\documentclass[11pt]{article}

\usepackage[T1]{fontenc}
\usepackage[utf8]{inputenc}
\usepackage[english]{babel}
\usepackage[margin=1in]{geometry}
\usepackage{lmodern}
\usepackage{microtype}
\usepackage{mathtools}
\usepackage{amsmath,amssymb,amsthm}
\usepackage{array}
\usepackage{booktabs}
\usepackage{tabularx}
\usepackage{float}
\usepackage{graphicx}
\usepackage{pdflscape}
\usepackage{caption}
\usepackage{enumitem}
\usepackage{fancyhdr}
\usepackage[numbers]{natbib}
\usepackage{doi}
\usepackage{xurl}
\usepackage{hyperref}
\usepackage{cleveref}
\usepackage{orcidlink}
\usepackage{titling}

% ---------------------------------------------------------------------------
% Series/cluster name macros (mirrors Foundations IV's \axisX / \chapX pattern)
% ---------------------------------------------------------------------------
\newcommand{\clusterA}{Temporal Currencies}
\newcommand{\clusterB}{Moving and Self-Generated Obstructions}
\newcommand{\clusterC}{Defect Integration and Emergent Transport}
\newcommand{\clusterD}{Productive Obstructions and Sparse Saturation}

\newcommand{\FVI}{Foundations~VI}
\newcommand{\FoundI}{Foundations~I}
\newcommand{\FoundII}{Foundations~II}
\newcommand{\FoundIII}{Foundations~III}
\newcommand{\FoundIV}{Foundations~IV}
\newcommand{\FoundV}{Foundations~V}
\newcommand{\NumLaws}{13}
\newcommand{\NumClusters}{4}

% Per-law short-name macros -- fill in / extend as chapters are drafted.
\newcommand{\lawGOne}{Hidden Amortized Solvency}
\newcommand{\lawGTwo}{Transfinite Escrow}
\newcommand{\lawGThree}{Amortized Potential Currency}
\newcommand{\lawGFour}{Moving-Cover Exhaustion}
\newcommand{\lawGFive}{Carry-Horizon Confinement}
\newcommand{\lawGSix}{Endogenous Needle Generation}
\newcommand{\lawGSeven}{Adversarial Mobility Confinement}
\newcommand{\lawGEight}{Odometer Abelianization}
\newcommand{\lawGNine}{Defect-Evacuation-to-Transport}
\newcommand{\lawGTen}{Lossful Boundary Persistence}
\newcommand{\lawGEleven}{Local-Rule Global-Anti-Symmetry}
\newcommand{\lawGTwelve}{Finite Witness Radiation}
\newcommand{\lawGThirteen}{Thin-Orbit Saturation}

% Sibling-paper macros (populate once/if sibling papers importing G-laws exist)
\newcommand{\PaperAOR}{the audited-operational-realisability paper}

\newtheorem{theorem}{Theorem}[section]
\newtheorem{lemma}[theorem]{Lemma}
\newtheorem{proposition}[theorem]{Proposition}
\newtheorem{corollary}[theorem]{Corollary}
\theoremstyle{definition}
\newtheorem{definition}[theorem]{Definition}
\newtheorem{remark}[theorem]{Remark}
\floatstyle{ruled}
\newfloat{algorithm}{tbp}{loa}
\floatname{algorithm}{Algorithm}

\urlstyle{same}
\captionsetup{
  font=small,
  labelfont=bf,
  labelsep=period
}
\captionsetup[table]{skip=6pt,position=top}
\captionsetup[figure]{skip=6pt}
\captionsetup[algorithm]{skip=6pt,position=top}
\crefname{algorithm}{algorithm}{algorithms}
\Crefname{algorithm}{Algorithm}{Algorithms}
\crefname{theorem}{theorem}{theorems}
\Crefname{theorem}{Theorem}{Theorems}
\crefname{lemma}{lemma}{lemmas}
\Crefname{lemma}{Lemma}{Lemmas}
\crefname{proposition}{proposition}{propositions}
\Crefname{proposition}{Proposition}{Propositions}
\crefname{corollary}{corollary}{corollaries}
\Crefname{corollary}{Corollary}{Corollaries}
\crefname{definition}{definition}{definitions}
\Crefname{definition}{Definition}{Definitions}
\crefname{remark}{remark}{remarks}
\Crefname{remark}{Remark}{Remarks}
\setlist[enumerate]{leftmargin=*, itemsep=2pt, topsep=4pt}
\setlength{\droptitle}{-3.2em}
\pretitle{\begin{center}\LARGE\bfseries}
\posttitle{\par\end{center}\vspace{0.5em}}
\preauthor{\begin{center}\normalsize}
\postauthor{\par\end{center}\vspace{-0.4em}}
\predate{\begin{center}\small}
\postdate{\par\end{center}\vspace{0.6em}}
\setlength{\emergencystretch}{2em}
\clubpenalty=10000
\widowpenalty=10000
\displaywidowpenalty=10000
\interfootnotelinepenalty=10000
\allowdisplaybreaks[2]
\fancypagestyle{firstpage}{\fancyhf{}
  % DOI is a deliberate pre-mint placeholder (design/08 SS3): filled at Zenodo mint, outside the
  % draft-completeness claim.
  \fancyfoot[C]{\scriptsize Automorph Inc., Wilmington, DE, USA \quad \texttt{ioannis@automorph.io} \quad \textcopyright\ 2026 \quad CC-BY 4.0 \quad \doi{10.5281/zenodo.PLACEHOLDER}}
  \renewcommand{\headrulewidth}{0pt}
  \renewcommand{\footrulewidth}{0pt}
}

\title{Six Birds Foundations VI:\\
A Catalog of Dynamical Structural Laws}
\author{Ioannis Tsiokos\,\orcidlink{0009-0009-7659-5964}}
\date{\today}
\hypersetup{
  pdftitle={Six Birds Foundations VI: A Catalog of Dynamical Structural Laws},
  pdfauthor={Ioannis Tsiokos},
  hidelinks
}

\begin{document}

\sloppy
\hbadness=10000

\maketitle
\thispagestyle{firstpage}

\begin{abstract}
The Six Birds corpus through \FoundIV{} certifies \emph{states}: each of its laws fixes a single map,
quotient, package, or declared finite family and proves a static classification about it. Nothing in
the corpus certifies \emph{runs} --- liveness (every orbit eventually earns a certificate), amortized
or transfinite cost over unbounded time, covers that must move with their targets, obstructions that
grow, react, or are manufactured by the system's own history, route defects that integrate into
potentials, transport that emerges from substrate dynamics, or content that persists through repeated
loss. This paper is the catalog of exactly those laws: \NumLaws{} dynamical structural laws
(G1--G13), organized in \NumClusters{} clusters, each stated in the six-part normal form inherited
verbatim from \FoundIV{} and each carrying a temporal quantifier that no static row possesses. The
laws were mined from eccentric mathematical problems --- Collatz, Goodstein, Egyptian fractions,
reverse-and-add, Recam\'an, the angel problem, sandpiles, rule 184, Ducci, aperiodic tilings,
Hadwiger--Nelson, Apollonian packings --- treated as theorem detectors: each problem's folklore
almost-proof names a promoted object, the object's lawful descended certificate fixes a boundary
expression, and the residual surviving comparison against every existing law is the law seed. All
thirteen laws are mechanized in Mathlib-free Lean~4 with per-declaration kernel-axiom audits; six
laws are full \texttt{theorem}-grade rows and one splits into an unconditional theorem and a schema,
four are \texttt{schema}-grade with named boxed hypotheses, and three are calibration-anchored
schemas discharged end-to-end by solved classical instances. The paper solves no open problem and
claims none: Collatz, Erd\H{o}s--Straus, Lychrel, Recam\'an, the Langton's-ant highway, Gilbreath,
and the exact chromatic number of the plane all remain open, and each law's nonclaims say so
explicitly.
\end{abstract}

\medskip
\noindent\textbf{Keywords.} dynamical structural law; liveness; amortized currency; moving-cover
exhaustion; adversarial confinement; endogenous obstruction; defect-to-transport certification; finite
witness radiation; layer-agnostic normal form; Six Birds theory.

% =============================================================================
\section{Introduction: the dynamical turn}\label{sec:intro}

A structural law, in the sense this series gives the term, is a layer-agnostic theorem: it quantifies
over declared abstract data --- carriers, quotients, ledgers, patch systems --- and classifies them,
so that any substrate supplying the data inherits the classification. \FoundIV{} assembled fifty-two
such laws \citep{Tsiokos2026FoundIV}, and one fact about that catalog, verified against its literal
text, is the starting point of this paper: every one of its Setups names fixed, singular data ---
one map, one quotient pair, one stability package, one \emph{declared finite} patch family --- and
every one of its theorems is a static biconditional or classification about that fixed data. No row
contains a step-indexed quantifier, a growing family, an adversary, a rate comparison, or an
emergence-over-time claim. The corpus certifies what a configuration \emph{is}; it is silent on what
a run \emph{does}.

That silence is not a small gap at the margin. It excludes the questions that animate whole regions
of mathematics and its applications: whether every orbit of a map eventually earns a descent
certificate; whether unbounded work can be paid for by a potential that is allowed to fluctuate;
whether termination can be audited when the very presentation of the state changes at every stage;
whether a cover can succeed when no finite package of patches suffices; whether prediction can
survive a horizon that recedes as the run advances; what happens when the obstruction to progress is
manufactured by the run's own history, or chosen adaptively by an adversary; when route disagreements
integrate into a canonical potential; when local disorder resolves into emergent transport; what
survives iterated loss; and how local rules, finite witnesses, and thin generators force, certify,
and saturate global structure. This paper records the thirteen laws --- the G-series --- that govern
exactly these questions, in four clusters: \clusterA{} (\Cref{sec:clusterA}), \clusterB{}
(\Cref{sec:clusterB}), \clusterC{} (\Cref{sec:clusterC}), and \clusterD{} (\Cref{sec:clusterD}).
\Cref{tab:whitespace} maps each law to the static row it most resembles and the dynamical content
that separates them; the per-law sections make each separation precise, and the rebadge discipline
behind the claim of newness is part of the method (\Cref{sec:methodology}).

\begin{table}[htbp]
\caption{The whitespace map: what each G-law adds that no static catalog row contains. Cousin rows
refer to \FoundIV's catalog \citep{Tsiokos2026FoundIV}.}
\label{tab:whitespace}
\footnotesize
\begin{tabularx}{\textwidth}{@{}lXX@{}}
\toprule
Law & Dynamical phenomenon & Closest static content, and the gap \\
\midrule
G1 & Liveness: every nonterminal orbit eventually solvent & F13a fixes one hiddenness interface at
one instant; no row has an unbounded eventual-solvency quantifier \\
G2 & Termination audited through changing presentations & F2 repairs one fixed quotient; the escrow
audit is stage-indexed and composed \\
G3 & Amortized cost over finite prefixes & F27 characterizes an orbit-constant conserved
quantity; here the potential is not orbit-constant and its differences telescope under audit \\
G4 & Exhaustion by covers that move with the target & F4 classifies one declared finite patch
system; here the family grows with the target \\
G5 & Prediction against a receding carry horizon & F13a/F3 fix one quotient pair or one route
comparison; here the failure family is step-indexed \\
G6 & Obstructions manufactured by the run's own history & F6 declares its needles up front; here the
needle is generated by the orbit and audited externally \\
G7 & Obstructions chosen adaptively by a budgeted adversary & F6 declares its stability package
up front with no turn index or reacting adversary; here certification is strategy-valued over
unbounded play \\
G8 & Route data integrating into a conserved odometer & F3 compares two routes once; F27 presumes
the conserved quantity; here the ledger is constructed \\
G9 & Local defects evacuating into emergent transport & F19 checks persistence of an
already-declared transport; here the transporter emerges from the run \\
G10 & Boundary content persisting through iterated loss & F34 legitimizes a single lossy step; here
one relationship survives every iterate via a regenerating invariant \\
G11 & Local rules forcing global aperiodicity & F4/F40 test one gluing or one symmetry descent; here
every candidate period receives its own finite defect \\
G12 & A finite obstruction's survival as a global certificate & F6 dissolves, prices, or confines
needles; here the needle's survival \emph{is} the lower-bound certificate \\
G13 & Thin generators saturating thick targets & F4 tests one local datum; here the claim is
asymptotic orbit coverage with typed residuals \\
\bottomrule
\end{tabularx}
\end{table}

The discovery instrument is as much a part of the contribution as the catalog. Each law was mined
from an eccentric mathematical problem --- a problem with a tiny rule and a global, temporal
obstruction --- by a fixed five-step protocol (\Cref{sec:methodology}): name the folklore almost-proof
that everyone reaches for and that fails; name the promoted object it secretly uses; state the lawful
descended certificate; isolate the residual that survives every existing law; and find the siblings,
at least three independent domains with at least one solved. The protocol is a filter, not a funnel
--- two candidate seeds died at the residual step and are recorded with their rejection reasoning
(\Cref{sec:methodology,sec:sketches}) --- and every cluster is anchored by solved, recognized
mathematics: Goodstein and Kirby--Paris for Cluster A's escrow, Dhar and Fey--Levine--Peres for
Cluster C's odometer, the fully settled angel problem for Cluster B's adversary, Berger through the
2023 monotiles and de~Grey's graph and Bourgain--Kontorovich for Cluster D. The open problems in the
catalog --- Collatz, Erd\H{o}s--Straus, Lychrel, Recam\'an, the ant's highway, Gilbreath, the exact
chromatic number of the plane --- appear only as conjectural instances whose hypotheses are named and
unverified; this paper proves nothing about any of them, and \Cref{sec:nonclaims} says so law by law.

A word on the paper's position in its series is required, because the nearest sibling is also the
easiest to confuse. Foundations~V (\emph{Endogenous Closure}) is \emph{carrier-side}: it states the
lawful conditions under which a cognitive, economic, or social system hosts its own closure apparatus
--- carries, prices, audits, and re-instates its own repair records \citep{FoundationsVCognition}.
This paper is \emph{theorist-side}: an external analyst certifies structural properties of a run ---
liveness, amortized cost, coverage, confinement, aperiodicity, radiation --- and at no point does any
G-law credit a system with certifying itself. The two projects fill different verified gaps (V:\ who
performs the closure operations; VI:\ nothing certifies runs), share no theorem, carrier
construction, or Lean namespace, and keep disjoint vocabularies; the one law whose name risks
adjacency, G6, carries an explicit disambiguation in its own section, and the two deliberate
cross-citation points (V's bounded-reflexivity law at the parked look-and-say sketch, V's curiosity
law at G6's exploration instantiation) are marked where they occur.

Everything stated as mechanized is mechanized: all thirteen laws are Lean~4 theorems or schemas in
the paper's companion repository, Mathlib-free with zero imports, with per-declaration kernel-axiom
audits (\Cref{app:formalization}); only two laws use classical choice, and the appendix records
which and why. Every numerical exhibit quoted in a law section is an exact-arithmetic lab run with a
committed prediction and seed, and the two-layer review record --- an internal gate on every
artifact, then an external adversarial cycle that found and closed two genuine gaps --- is reported
in \Cref{subsec:reviewer-cycle-log} rather than smoothed away. The reader who wants the apparatus
first should read \Cref{sec:apparatus,sec:methodology} and then any cluster; the reader who wants
the mathematics first can start at \Cref{sec:clusterC} (the solved-anchor cluster) and consult the
apparatus as needed; \Cref{sec:currency-taxonomy} states the one genuinely new theorem outside the
catalog proper, and \Cref{sec:instantiations-glance} is the navigational index over all sixty-five
instantiations.

% =============================================================================
\section{Apparatus and conventions}\label{sec:apparatus}

This section fixes the Six Birds vocabulary the catalog uses, at the minimum depth needed to read the
laws; the defining papers remain authoritative. A \emph{theory package} is a tuple
$\mathcal{T} = (Z, f, \Sigma_f, E, \mathcal{A})$: a finite carrier of microstates, a deterministic
lens, the definability algebra of the lens, an idempotent completion/packaging endomap whose fixed
points are the layer's objects, and an audit functional monotone under coarsening
\citep{Tsiokos2026EmergenceCalculus}. Six primitive structural roles recur across layers: \textbf{P1}
update-versus-quotient compatibility (descent obstruction); \textbf{P2} macro-specification
representability (constraints, gating, feasibility); \textbf{P3} enriched route/coherence mismatch
(holonomy) --- diagnostic only, explicitly not a directionality certificate; \textbf{P4}
stage/order/refinement structure; \textbf{P5} quotient packaging; \textbf{P6}
audit/bookkeeping/provenance, whose affinity-carrying, detailed-balance-breaking face
(\emph{P6\_drive}) is reserved strictly for thermodynamic-arrow content and never inferred from P3
data \citep{Tsiokos2026FoundII}.\footnote{The primitive numbering drifted between the corpus's
earliest paper and the later foundations series; this paper uses Foundations~II's corrected meanings
throughout, in particular for P3 and P4.}

\paragraph{Quotients and hiddenness.} On a history space, the \emph{current quotient} identifies
histories indistinguishable by every declared current observable; the \emph{predictive quotient}
identifies histories indistinguishable under all admissible future continuations, and always refines
the current one. A \emph{predictive witness} --- two histories equal currently but differing
predictively --- is exactly what it means for hidden predictive content to exist; the load-bearing
predicate is binary exposure, not coupling strength \citep{Tsiokos2026Holonomy,Tsiokos2026Hiddenness}.
An update $F$ \emph{descends} through a quotient $q$ exactly when a quotient-level update $\bar F$
satisfies $q \circ F = \bar F \circ q$, equivalently when the split-pair obstruction
$\Delta(q, F, r) = \{(x, x') : q(x) = q(x') \wedge r F(x) \neq r F(x')\}$ is empty
\citep{Tsiokos2026FoundIV}. Several G-laws are step-indexed extensions of exactly this test: G1 says
no fixed residue quotient is future-sufficient for Collatz descent probes, and G5's carry horizon is
a step-indexed \emph{sequence} of quotient-prediction failures, where the corpus fixes one pair.

\paragraph{Promotion, SAU, and the boundary certificate.} \emph{Promotion} is passage to a richer
layer --- $\mathbb{C}$ from $\mathbb{R}$, ordinals from $\mathbb{N}$, the $2$-adics from the odd
integers --- via a recorded bridge. The \emph{SAU pattern} (Strict Audited Utility): a promoted
object $U$ that does not descend may still legally solve a lower-layer problem, provided a selected
boundary expression $C(U)$ \emph{does} descend to a lower-layer answer; the audit never licenses
treating $U$ itself as a lower-layer object \citep{Tsiokos2026WhyMath}. Throughout this paper we call
that descended, natively checkable assertion the \textbf{boundary certificate} --- the term is this
paper's own, introduced here once (the F-series papers use per-role certificate vocabulary instead),
and it is used in a fixed sense: the native inequality, equality, or finite record that promoted
support is allowed to purchase. Every Cluster A law is a SAU specialization in which the promoted
object is \emph{temporal} --- a future-valuation ledger, an ordinal escrow, a potential --- and the
boundary certificate is a descent or termination statement; the $2$-adic ghost $-1$ of
\Cref{sec:clusterA:g1} is the canonical non-descending promoted object.

\paragraph{Currencies.} A \emph{currency} is a low-dimensional map $u : Y \to \mathbb{R}^k$ playing
at least one of four structural roles --- C1 constraint, C2 ledger, C3 potential, C4 dual/shadow
price; statistics that merely correlate with performance are proxies, not currencies
\citep{Tsiokos2026Currency}. The corpus's currencies are all real-valued and step-monotone or
budget-bounded; \Cref{sec:currency-taxonomy} records the three temporal roles (C5 amortized-potential,
C6 transfinite-escrow, C7 solvency) that Cluster A forces.

\paragraph{Needles, membranes, and lawful repairs.} A \emph{needle} is a localized obstruction
surviving all native-layer probes; the corpus discipline prices needles against a declared budget,
dissolves them by promotion, or confines them by duality \citep{Tsiokos2026NeedleKiller}. A
\emph{membrane} is a budget condition holding at every accepted stage, and its failure is a critical
pair with exactly five lawful repairs (declared nonclaim, raised budget, rewritten carrier/audit,
genuinely new probes, completed transport family) --- same-package repetition is not a repair. In the
corpus every needle, budget, and probe family is declared up front and static; Clusters B and D live
precisely in the gap this leaves (history-generated needles, adversarial needles, and needles whose
survival is a positive certificate).

\paragraph{The catalog form, status tags, and discharge vocabulary.} \FoundIV{} presents its
fifty-two laws in a six-part normal form --- Setup, Theorem, Proof spine, Case enumeration,
Interaction status, Nonclaims --- each with five worked layer instantiations
\citep{Tsiokos2026FoundIV}; this paper adopts the template verbatim. One verified fact about that
catalog does structural work everywhere here: every F-IV Setup names fixed, singular data --- one
map, one quotient pair, one declared finite patch system --- and every F-IV theorem is a static
classification of that fixed data; no row contains a step-indexed quantifier, a growing family, an
adversary, a rate comparison, or an emergence-over-time claim. The temporal quantifier is therefore
the novelty axis of every G-law, and each law's Setup states it explicitly. Status tags are
three-valued and typographically fixed: \texttt{theorem} (the abstract implication is proved),
\texttt{schema} (named hypotheses carry the implication; no concrete instance is asserted to satisfy
them until discharged), and \texttt{calibration-anchored schema} (a schema whose mechanism is fully
discharged by a recognized solved classical instance). Residual accounting uses the AOR status
vocabulary \citep{Tsiokos2026AOR}, including the four asymptotic statuses recalled in
\Cref{sec:currency-taxonomy}, and the corpus nonclaim register is in force throughout --- most
prominently: route mismatch never implies drive content (G8's bridge is conditional on its named
hypotheses, which is legal); local admissibility never implies global admissibility without
exhaustion data (G4, G11); no same-level self-certification (G6, G7 are theorist-side); and promoted
support is never literalized at the native layer (all of Cluster A)
\citep{Tsiokos2026FoundIII}.

% =============================================================================
\section{Methodology: eccentric problems as theorem detectors}\label{sec:methodology}

This section has no counterpart in \FoundIV{}, and the addition is deliberate: the catalog's laws were
not invented and then illustrated, they were \emph{mined}, and the mining procedure is a first-class
contribution of this paper. The raw material is the class of eccentric problems --- Collatz,
Goodstein, Erd\H{o}s--Straus, Lychrel, Recam\'an, the angel problem, sandpiles, Langton's ant, Ducci,
Gilbreath, aperiodic tilings, Hadwiger--Nelson, Apollonian packings --- problems that are easy to
state, resistant in a characteristic way, and each equipped with a folklore argument that everyone
reaches for and that fails. The thesis of the method is that the folklore argument's failure is
structured: what it illegitimately assumes is a \emph{promoted object}, and what would repair it is a
\emph{descended certificate}, and the gap between the two, once it survives comparison against every
existing law of the corpus, is a new law.

For each problem, in order:
\begin{enumerate}
  \item \textbf{Identify the tempting almost-proof} --- the argument everyone reaches for that fails.
    The almost-proof reveals the illegal overread.
  \item \textbf{Name the promoted object the almost-proof secretly uses} --- an invariant measure on a
    completion, a fixed finite certificate package, a well-founded escrow. Typically the object does
    not descend to the native carrier.
  \item \textbf{State what a lawful descended certificate would be.} This fixes the boundary
    expression: the native, checkable assertion that promoted support is allowed to purchase
    (\Cref{sec:apparatus}).
  \item \textbf{Isolate the residual that survives every existing law.} Run the candidate against the
    full Foundations~IV catalog (exact statements, not summaries), the Foundations~V laws, the
    nonclaim register, and the AOR status list. What remains unhandled is the law seed --- and this
    step kills candidates as often as it produces them.
  \item \textbf{Find the siblings.} The seed is a law only if the same residual shape appears in at
    least three independent domains, at least one of them solved (the calibration instance).
    Otherwise it is a substrate observation, not a law.
\end{enumerate}
Admission then passes four further gates inherited from the corpus's discipline: a precise recognition
source making every Setup datum definite; a rebadge test requiring the candidate to differ from every
existing law on at least one declared axis --- for all thirteen laws here the differentiating axis is
the same, a step-indexed, growing, adversarial, or rate-comparison witness schema that no static
catalog row possesses, which is why an editing pass that accidentally deletes a temporal quantifier
from a law's statement collapses the law into its cousin and fails the gate; nonclaim compatibility
(a law that needs an unconditional version of a forbidden inference is dead); and a replication
ladder on the instantiation table (at least one solved calibration plus at least one instance with
genuine engagement beyond citation). The remainder of this section runs the protocol on four seeds:
two that became laws, and two that the protocol correctly killed.

\subsection*{Worked example: Collatz to G1}

\emph{Step 1.} The tempting almost-proof is negative drift. Over the accelerated odd map
$T(n) = (3n+1)/2^{a(n)}$, the valuation $a$ behaves on average like a geometric variable with mean
$2$, so the expected step effect on $\log n$ is $\log 3 - 2 \log 2 = \log \tfrac34 < 0$: orbits
``should'' shrink. The argument fails pointwise --- expectation controls almost-all behavior at best,
and the conjecture is a pointwise liveness claim.

\emph{Step 2.} The promoted object is the invariant measure on a completion: the drift computation
lives on $\mathbb{Z}_2$, where the valuation statistics are exact, not on the odd integers. On that
completion the dynamics genuinely have non-descending points --- the ghost $-1$ of
\Cref{sec:clusterA:g1} is a fixed point --- so the measure-level argument cannot see the pointwise
obstruction it is being asked to settle.

\emph{Step 3.} The lawful descended certificate is the solvency inequality: a finite $k$ with
$V_k(n) > 0$, equivalently the exact integer comparison $2^{A_k} n > 3^k n + B_k$, checkable at the
native layer with no reference to the completion.

\emph{Step 4.} The residual against the existing corpus is the unbounded eventual-solvency
quantifier. The closest static law fixes one hiddenness interface at one instant; nothing in the
corpus asks whether a time-indexed family of hidden-debt membranes eventually breaks for every
nonterminal point. That residual --- membranes $N_k$, ghost locus, native/ghost separation --- is
G1's content.

\emph{Step 5.} The siblings: amortized analysis (binary counters --- solved), size-change
termination (solved), aliquot dynamics (open, with a structurally richer ledger), switched-systems
certificates (structural). Three-plus domains, one solved: G1 is admitted, with its Part~(a)
reduction unconditional and its Part~(b) discharge honestly split off as a schema.

\subsection*{Worked example: sandpiles to G8, from the solved end}

The protocol also runs on solved problems, where it extracts the law rather than the obstruction.
For abelian sandpiles the tempting almost-proof is ``the firing order obviously cannot matter'' ---
empirically true and non-explanatory (step 1). The promoted object is the global stabilization
outcome itself, presumed order-independent before any hypothesis is named (step 2). The descended
certificate is the odometer: the per-site firing count vector, checkable on any finite run
(step 3). The residual against the corpus (step 4) is route \emph{integration} over an unbounded
run: the static catalog can compare two routes once and record a residue, or characterize an
already-given conserved quantity, but nothing constructs a canonical ledger from all legal routes of
an unbounded run --- that requires the abelian compatibility hypothesis, and isolating exactly that
hypothesis (with the case-(b) witness showing confluence alone does not suffice) is what turned a
solved curiosity into \Cref{thm:G8}. The siblings (step 5): rotor-routing (solved, same family),
commuting-exchange subfamilies of load balancing, settlement netting --- with Dhar's theorem
supplying the calibration, so the law lands at \texttt{theorem} grade with solved mathematics at its
core.

\subsection*{Negative walkthrough: Hidden-Atlas Exactness (absorbed)}

The protocol's value shows most clearly on seeds it kills. Somos-sequence integrality is an
``integrality miracle'': recurrences that divide at every step yet provably stay integral. Step 1:
the almost-proof is ``the divisions just happen to work out'' --- observational, no mechanism.
Step 2: the promoted object is a hidden atlas of cluster-algebra charts, under which each term is
manifestly a Laurent polynomial in the seed variables --- integrality is a shadow of Laurentness, not
a numerical accident. Step 3: the descended certificate is the Laurent-positivity certificate at
each mutation step, descending to native integrality once denominators formally cancel. Step 4 is
where the candidate dies: the corpus already owns this shape --- a promoted object that does not
descend but whose selected boundary expression does is precisely the SAU pattern of
\Cref{sec:apparatus}, and ``promoted mutable chart, hidden cancellation, native integrality'' is that
pattern specialized to cluster algebras, with no residual left over. Verdict: not a law; a case study
for the SAU registry. The mining discipline is seed-killing as much as seed-generation, and a
candidate that namechecks rich mathematics (cluster algebras, integrable systems) is not thereby a
law.

\subsection*{Negative walkthrough: Self-Observation Atomization (parked)}

The second kill has a different shape. Conway's look-and-say sequence eventually decomposes into a
fixed finite chemistry of elementary atomic subsequences --- the Cosmological Theorem. Step 1: the
almost-proof is ``self-description obviously stabilizes eventually.'' Step 2: the promoted object is
qualitatively new --- the system's own rule is semantic, the object functioning as a theory of
itself; this is genuine self-reference, not a passive auxiliary structure. Step 3: the descended
certificate would be a bounded parsing grammar under which every long dependency splits, stabilizes
into a finite atom, or is forbidden. Step 4 kills the candidate by \emph{collision} rather than
absorption: the residual lands squarely in Foundations~V's territory of bounded reflexivity ---
carrier-hosted audit structures certifying their own soundness \citep{FoundationsVCognition} ---
and brushes the corpus's carrier-side circularity limits; disentangling V's carrier-side
self-reference apparatus from VI's theorist-side certification apparatus is real research, not
something the mining protocol resolves inline. Verdict: parked as a deferred sketch with an explicit
pointer to Foundations~V (\Cref{sec:sketches}), exactly as \FoundIV{} parks its own boundary
candidates --- deferred, not dismissed. The two negative walkthroughs exhibit the protocol's two
distinct failure modes: \emph{absorbed} (already owned by existing apparatus) and \emph{collided}
(contested between live research programs).

% =============================================================================
\section{Cluster A --- \clusterA}\label{sec:clusterA}

The corpus is rich in currencies --- constraints, ledgers, potentials, duals --- but every one of them
is checked at an instant or across one declared comparison. What it lacks are the \emph{temporal}
currencies: the solvency question (does an unbounded run eventually earn a native descent
certificate?), the escrow question (can a changing presentation be audited by strict descent in a
well-founded carrier?), and the credit question (when does a supplied potential bound the total cost of
every finite prefix?). The three laws of this cluster record those currencies. \lawGOne{} (G1) is the
liveness law and the flagship of the catalog: it splits the eventual-solvency question into an
unconditional reduction and a conditional discharge schema whose named hypotheses are exactly what
remains open for Collatz. \lawGTwo{} (G2) is the termination law for stage-indexed runs, with
Goodstein sequences discharging the full transfinite mechanism. \lawGThree{} (G3) is the safety law:
the classical potential-method telescoping identity, stated as a currency with an explicit legality
gate.

The grades are again mixed by design. G2 and G3 carry status tag \texttt{theorem}. G1 splits: its
Part~(a) reduction is an unconditional \texttt{theorem}, while its Part~(b) discharge is a
\texttt{schema} --- calibrated in mechanism by solved amortized-analysis instances and by a finite
ghost toy, but not yet by the real Collatz discharge, whose two boxed hypotheses remain open.

\subsection{\lawGOne}\label{sec:clusterA:g1}

Let $X$ be a native discrete carrier with update $T : X \to X$ and an accepted terminal set
$A \subseteq X$ fixed by $T$. A run from $x$ is $x_0 = x$, $x_{j+1} = T(x_j)$. The promoted ledger
records per-step visible debt $d_j$, hidden payment $p_j$, and a finite-prefix tax term
$\operatorname{tax}_k(x)$; the \emph{solvency currency} is
\[
  V_k(x) \;=\; \sum_{j<k} (p_j - d_j) \;-\; \operatorname{tax}_k(x),
\]
and a descent certificate is valid only when the boundary inequality is exact:
\[
  V_k(x) > 0 \quad\Longleftrightarrow\quad T^{k}(x) < x
\]
in the declared native order. The unbounded quantifier is the law's content: G1 asks whether every
nonterminal native $x$ eventually has some finite $k$ with $V_k(x) > 0$. For each $k$, the
\emph{bad-tail membrane} is
\[
  N_k \;=\; \{\, x \in X \setminus A \;:\; V_j(x) \le 0 \ \text{for every } 1 \le j \le k \,\},
\]
so $x \in \bigcap_k N_k$ says that a nonterminal point never becomes solvent at any finite stage.
Finally, let $\widehat{X}$ be a completion containing the limits of infinite insolvent threads, and
let the \emph{ghost locus} $\Gamma$ be the non-descending fixed or periodic points of
$\widehat{X} \setminus X$.

For the Collatz flagship, $X$ is the odd positive integers, $A = \{1\}$, and the accelerated map
\[
  T(n) \;=\; \frac{3n+1}{2^{a(n)}}, \qquad a(n) = \nu_2(3n+1),
\]
is closed on $X$ with $T(1) = 1$, so the known trivial cycle is a genuine fixed point, terminal before
any membrane test. Writing $n_j = T^{j}(n)$, $a_j = \nu_2(3 n_j + 1)$, $A_k = \sum_{j<k} a_j$, and
$B_0 = 0$, $B_{j+1} = 3 B_j + 2^{A_j}$, an induction gives the exact affine ledger
\[
  n_k \;=\; \frac{3^{k} n + B_k}{2^{A_k}},
\]
and with $V_k(n) = A_k \log 2 - k \log 3 - \log\!\bigl(1 + B_k / (3^k n)\bigr)$ one computes
$\log(n_k / n) = -V_k(n)$, so
\[
  V_k(n) > 0
  \quad\Longleftrightarrow\quad
  n_k < n
  \quad\Longleftrightarrow\quad
  2^{A_k} n \;>\; 3^{k} n + B_k .
\]
The logarithmic form is presentation; the certificate checked in the lab is the exact integer
inequality on the right. The completion is $\widehat{X} = \mathbb{Z}_2$, where $-1 = \ldots 1111_2$ is
a ghost: $3(-1) + 1 = -2$ has $\nu_2 = 1$, so $T(-1) = -1$, a non-descending fixed point outside the
native carrier.

The closest Foundations~IV cousin is F13a (Hiddenness Normal Form), which fixes one instant: a current
quotient, a predictive refinement, and predicates asking whether a current fiber mixes predictive
classes. G1 is that hiddenness interface stretched over an unbounded run --- the hidden data are the
future valuation ledger, the membranes are indexed by time, and the question is whether hidden debt is
eventually discharged by a native certificate. The Collatz instance makes the hiddenness concrete:
computing $T(n) \bmod 2^m$ requires knowing $n \bmod 2^{m + a}$, and the valuation $a$ is unbounded
over odd inputs, so no fixed residue quotient $n \bmod 2^m$ is future-sufficient for all descent
probes. Part~(a) below carries status tag \texttt{theorem}; Part~(b) carries \texttt{schema}.

\begin{theorem}[\lawGOne, Part (a): bad-tail reduction]\label{thm:G1}
Pointwise liveness is equivalent to emptiness of the infinite bad-tail intersection:
\[
  \bigl(\;\forall\, x \in X \setminus A ,\ \exists\, k \ge 1 :\ V_k(x) > 0\;\bigr)
  \quad\Longleftrightarrow\quad
  \bigcap_{k \ge 1} N_k \;=\; \varnothing .
\]
Points of $A$ are terminally discharged before the membrane test.
\end{theorem}

\begin{proof}
If liveness holds and some nonterminal $x$ lay in every $N_k$, the finite $k$ with $V_k(x) > 0$
contradicts the membership condition of $N_k$ itself. Conversely, if the intersection is empty, each
nonterminal $x$ fails some $N_k$, and unfolding that failure produces $1 \le j \le k$ with
$V_j(x) > 0$ --- a finite descent certificate. The two directions use nothing beyond the membrane
definition; the reduction is unconditional.\footnote{Anchored to the substrate Lean at
\path|lean/SixBirdsFoundationsVI/Laws/G1HiddenAmortizedSolvency.lean|: \path|part_a_reduction|, over
the setup predicates \path|Descends|, \path|BadTailMembrane|, \path|PointwiseLiveness|,
\path|EmptyBadTailIntersection|. Status tag \texttt{theorem}.}
\end{proof}

\begin{theorem}[\lawGOne, Part (b): ghost-discharge schema]\label{thm:G1b}
Assume:
\begin{itemize}
  \item \textup{[H-G1-ghost-convergence]} for a computable neighborhood system $U_h(\Gamma)$ in
    $\widehat{X}$ with $\bigcap_h U_h(\Gamma) = \Gamma$: every infinite bad thread based at a
    nonterminal native point has a limit in $\widehat{X}$, every such limit lies in $\Gamma$, and
    hence for every $h$ the thread is eventually in $U_h(\Gamma)$;
  \item \textup{[H-G1-separation]} for every nonterminal native point $x$ there are computable
    $h(x)$ and $K(x)$, established by native certificates only, with
    \[
      T^{t}(x) \notin U_{h(x)}(\Gamma) \qquad \text{for every } t \ge K(x).
    \]
\end{itemize}
Then $\bigcap_k N_k = \varnothing$; by Part~(a), every nonterminal native point is live, and for each
fixed $x$ the bad-tail predicate $x \in N_k$ is eventually false. For Collatz both hypotheses are
unproved: they are named obligations, not evidence that the conjecture has been settled.
\end{theorem}

\begin{proof}
Suppose some nonterminal $x$ lies in every $N_k$; its orbit is an infinite bad thread. Ghost
convergence puts the thread eventually inside $U_{h(x)}(\Gamma)$; native separation keeps the same
orbit outside $U_{h(x)}(\Gamma)$ from time $K(x)$ on. A time beyond both horizons is a contradiction,
so no such $x$ exists, and Part~(a) converts the emptiness into liveness.\footnote{Anchored to the
substrate Lean at \path|lean/SixBirdsFoundationsVI/Laws/G1HiddenAmortizedSolvency.lean|:
\path|part_b_discharge|, with hypothesis predicates \path|GhostConvergence| and
\path|NativeSeparation|. Status tag \texttt{schema}.}
\end{proof}

Two unconditional Collatz facts calibrate the schema's difficulty from opposite directions. First, the
ghost is \emph{shadowed} arbitrarily long by native points: for $n_L = 2^L - 1$ one verifies
inductively that
\[
  n_i \;=\; 3^{i}\, 2^{\,L-i} - 1 \qquad (0 \le i < L),
\]
with $a_i = 1$ at every step $i < L - 1$, so the orbit spends at least $L - 1$ consecutive accelerated
steps insolvent --- the binary expansion $\underbrace{1\cdots1}_{L}$ mimics the $2$-adic ghost $-1$
for $L - 1$ steps. Arbitrarily long insolvent prefixes therefore exist; this is evidence against every
fixed finite-depth proof strategy, and evidence neither for nor against the conjecture itself. Second,
the affine ledger globalizes the cycle question: a valuation word $(a_0, \ldots, a_{k-1})$ can close
into a positive integer cycle only if
\[
  n \;=\; \frac{B_k}{2^{A_k} - 3^{k}}
\]
is a positive integer (in particular $2^{A_k} > 3^k$) --- a local-to-global integrality residual, not
a liveness statement.

The orbit fates partition by first discharge time:

\begin{center}
\footnotesize
\begin{tabularx}{\textwidth}{@{}lXX@{}}
\toprule
Case & Status & Witness type \\
\midrule
(a) Terminal or immediately solvent & $x \in A$, or $x \notin A$ with $V_1(x) > 0$, hence
$T(x) < x$. & Terminal membership, or a one-step ledger certificate. \\
(b) Finite bad prefix before discharge & Least discharge time $k > 1$: $x \in N_{k-1}$ and
$V_k(x) > 0$. & Bounded membrane audit plus the exact ledger pair $(A_k, B_k)$; the $2^L - 1$ family
shows $k - 1$ can be arbitrarily large. \\
(c) Infinite bad thread & $x \in \bigcap_k N_k$; under Part~(b)'s hypotheses, ghost-bound and then
excluded. & Infinite membrane record; for Collatz, exactly the unproved obstruction. \\
\bottomrule
\end{tabularx}
\end{center}

In interaction terms, G1 is the promoted-support pattern in temporal form: the ledger, the completion
$\mathbb{Z}_2$, and the ghost $-1$ are promoted objects that never descend to the native layer; the
only native assertion is the boundary certificate $V_k(x) > 0$, equivalently $T^k(x) < x$. The
computational record is deliberately split across four labs. G1-L1 checked all $500000$ odd
$n < 10^{6}$ ($499999$ nonterminal starts, $2241405$ exact ledger checkpoints): every nonterminal
start reached first descent, with maximum first-descent step $111$ at $n = 626331$. G1-L2 extended to
all $4999999$ nonterminal odd starts below $10^{7}$ (maximum $k^{*} = 155$ at $n = 8088063$) and
verified the $2^{L}-1$ shadowing formula for $L = 2, \ldots, 80$ --- while \emph{falsifying} the
over-strong prediction $k^{*} = L - 1$ in all $79$ tested cases, a recorded miss kept as a discipline
exhibit. G1-L3 exhaustively classified the $32768$ odd residues modulo $2^{16}$ in a finite ghost toy:
$32765$ residues descended, and the two exceptions $\Gamma_{16} = \{14563, 21845\}$ form the corrected
finite ghost locus, with no genuine nonterminal bad cycle --- a completed finite-model exhibit of the
discharge mechanism, not a calibration of real Collatz. G1-L4 probed the aliquot instantiation with
$100000$ starts: $82728$ terminated at $0$, $2085$ entered cycles ($22$ distinct cycles observed), and
$15187$ escaped the $10^{12}$ cap, matching the open Catalan--Dickson picture.

\subsection*{Layer instantiations}\label{layer-inst:G1}
\begin{description}
  \item[Collatz accelerated odd map --- conjectural flagship.] The affine ledger, the $2$-adic ghost
    $-1$, the $2^{L}-1$ shadowing family, and the cycle integrality condition are verified algebraic
    facts; liveness remains open because both boxed hypotheses are unproved
    \citep{Lagarias1985,Lagarias2010,Tao2019Collatz}.
  \item[Binary-counter increment --- calibration of the currency mechanism.] With $\Phi$ the number of
    $1$-bits, an increment across $r$ trailing ones costs $r + 1$, changes potential by $1 - r$, and
    is exactly solvent at every step; this calibrates the solvency currency, not the ghost mechanism
    \citep{Tarjan1985,CormenEtAl}.
  \item[Size-change and ranking-function termination --- calibration.] Termination analysis discharges
    local non-descent by a hidden well-founded ranking certificate --- ``promoted ledger gives native
    liveness'' as solved technology, without a Collatz-style ghost \citep{LeeJonesBenAmram2001}.
  \item[Aliquot dynamics --- conjectural folded instance.] For $s(n) = \sigma(n) - n$ the
    Catalan--Dickson question is open, and the natural currency is the whole factorization/divisor
    ledger rather than a scalar valuation --- G1 ledgers need not be one-dimensional
    \citep{GuySelfridge1975,GuyUnsolved}.
  \item[Switched and probabilistic termination certificates --- structural.] Average-dwell-time and
    ranking-supermartingale methods separate hidden average certificates from pointwise native
    progress; they become G1 instances only with an exact boundary certificate and pointwise residual
    \citep{Liberzon2003,AgrawalChatterjeeNovotny2017}.
\end{description}

The Collatz nonclaims (including the status of Tao's almost-all result as \texttt{probabilistic}
relative to G1's pointwise target), the finite-toy scope of G1-L3, and the promoted-object discipline
are consolidated in \Cref{sec:nonclaims}.

\subsection{\lawGTwo}\label{sec:clusterA:g2}

Let $X$ be a native state carrier with a stage-indexed run $x_0, x_1, x_2, \ldots$,
$x_{k+1} = T_k(x_k)$ (a fixed transition is the special case $T_k = T$), and accepted terminal set
$A \subseteq X$. For each stage $k$, declare a presentation $\pi_k : X \to P_k$ and an escrow map
$e_k : P_k \to W$ into a well-founded carrier $(W, <)$. The audit condition below is deliberately
imposed on the \emph{composed} object --- presentation change plus native step, checked as one
transition --- because a changing presentation can hide descent by moving value between encodings; a
proof that decreases only before or only after an unrecorded rebasing has not discharged the audit.
The temporal quantifier is unbounded: G2 asks whether every infinite nonterminal run would force an
infinite strictly decreasing chain in $W$. The Foundations~IV cousin F2 (Descent--Repair) fixes one
quotient and one map and repairs a single descent obstruction; G2 changes the logical type by making
the quotient stage-indexed and the descent condition run-length. The row carries status tag
\texttt{theorem}.

\begin{theorem}[\lawGTwo]\label{thm:G2}
Assume:
\begin{itemize}
  \item \textup{[H-G2-stage-coherence]} for every nonterminal step of the run,
    \[
      e_{k+1}\bigl(\pi_{k+1}(T_k(x_k))\bigr) \;<\; e_k\bigl(\pi_k(x_k)\bigr);
    \]
  \item $(W, <)$ is well-founded.
\end{itemize}
Then no infinite nonterminal run exists: every run reaches $A$ in finitely many steps. When $W$
carries an effective height extractor for the initial escrow value, the extractor yields a finite
step-count certificate or bound; otherwise the descended certificate is the finite terminating trace
itself.

For declared problem families, the existence of a sound uniform escrow may carry a named minimum order
strength. For Goodstein sequences that strength is $\varepsilon_0$: the standard proof descends
through ordinals cofinal below $\varepsilon_0$, and by Kirby--Paris, Goodstein's theorem is not
provable in Peano Arithmetic --- equivalently, for this normal form, no PA-internal escrow whose
well-foundedness is provable below the $\varepsilon_0$ strength can certify termination of all
Goodstein sequences. The price is imported proof theory, recorded here, not reproved.
\end{theorem}

\begin{proof}
If an infinite nonterminal run existed, stage coherence would string its escrow values into an
infinite strictly descending chain
\[
  e_0(\pi_0(x_0)) \;>\; e_1(\pi_1(x_1)) \;>\; e_2(\pi_2(x_2)) \;>\; \cdots
\]
in $W$, contradicting well-foundedness. No monotonicity of the native values is needed --- only the
audited escrow descends.

The Goodstein calibration discharges every ingredient. Take $X = \mathbb{N}$, $A = \{0\}$, stage base
$k \ge 2$; the presentation $\pi_k(n)$ is the hereditary base-$k$ expansion, the transition is
$G_k(n) = \operatorname{rebase}_{k \to k+1}(n) - 1$, and the escrow $O_k(n)$ replaces every base
symbol in the hereditary expansion by $\omega$. Rebasing preserves the ordinal,
\[
  O_{k+1}\bigl(\operatorname{rebase}_{k \to k+1}(n)\bigr) \;=\; O_k(n),
\]
because hereditary syntax is preserved position-by-position, while subtracting $1$ from a positive
hereditary expression strictly decreases the ordinal; hence
$O_{k+1}(G_k(n)) < O_k(n)$ --- exactly the composed audit. (At base $2$, starting from $4 = 2^{2}$:
$O_2(4) = \omega^{\omega}$; rebasing gives $3^{3} = 27$ with $O_3(27) = \omega^{\omega}$ unchanged,
and the subtraction $27 - 1 = 26 = 2 \cdot 3^2 + 2 \cdot 3 + 2$ drops the escrow to
$\omega^{2} \cdot 2 + \omega \cdot 2 + 2 < \omega^{\omega}$, while the native integer has jumped from
$4$ to $26$.) The escrow values are ordinals below $\varepsilon_0$, well-founded in the usual
set-theoretic sense, so every Goodstein sequence terminates at $0$; the PA-independence price is the
cited Kirby--Paris theorem \citep{KirbyParis1982}.\footnote{Anchored to the substrate Lean at
\path|lean/SixBirdsFoundationsVI/Laws/G2TransfiniteEscrow.lean|: \path|StageCoherence|,
\path|wellFounded_no_descending_chain|, \path|no_infinite_nonterminal_run|,
\path|not_nonterminal_of_stage_coherence|. Status tag \texttt{theorem}.}
\end{proof}

The escrow audits classify by the least order strength that actually carries them:

\begin{center}
\footnotesize
\begin{tabularx}{\textwidth}{@{}lXX@{}}
\toprule
Case & Order-strength status & Witness type \\
\midrule
(a) Natively monotone & $W$ embeds in $\mathbb{N}$. & Scalar ranking function with strict
decrease. \\
(b) Promoted scalar / lexicographic & $W$ embeds in a finite lexicographic product such as
$\mathbb{N}^m$, and no one-dimensional rank has been supplied. & Lexicographic or multicomponent
ranking certificate. \\
(c) Transfinite escrow & A genuine ordinal notation system is used (order type $\ge \omega^\omega$;
Goodstein is cofinal below $\varepsilon_0$). & Stage-coherent ordinal assignment such as $O_k$. \\
(d) No escrow at declared strength & The declared strength cannot soundly certify the whole family. &
Lower-bound or independence witness; Kirby--Paris for PA-strength escrows on Goodstein. \\
\bottomrule
\end{tabularx}
\end{center}

In interaction terms, G2 couples P4 staging/refinement to P1 descent; the escrow is P6 ledger content
whose values live in a well-founded order rather than a numeric budget, and the descended boundary
certificate is finite termination of the native run. The lab record: G2-L1 checked Goodstein starts
$1, \ldots, 30$ under a step budget of $4$, verifying all $116$ per-step rebase-invariance and
strict-descent identities exactly; starts $1$ and $2$ terminated within the budget, and the runs
reached base $6$ with a maximum integer of $120605$ bits (about $36306$ decimal digits). This is a
bounded identity exhibit, not a computational termination-length claim.

\subsection*{Layer instantiations}\label{layer-inst:G2}
\begin{description}
  \item[Goodstein sequences and hydra games --- calibration.] Hereditary-base presentation, ordinal
    escrow below $\varepsilon_0$, and the Kirby--Paris price give the complete positive-and-price
    mechanism \citep{KirbyParis1982}.
  \item[Ranking-function termination --- calibration.] Floyd-style arguments assign well-founded ranks
    to program states; ordinal-valued variants are G2 instances when the rank is stage-coherent across
    control-state presentations \citep{Floyd1967}.
  \item[Term-rewrite termination by path orderings --- calibration.] Recursive path orderings supply
    well-founded term measures of transfinite order type; a rewrite step is legal when the left side
    strictly dominates the right \citep{Dershowitz1982}.
  \item[Proof-theoretic ordinal analysis --- structural.] Gentzen's consistency analysis of Peano
    Arithmetic by induction up to $\varepsilon_0$ is the proof-theoretic background of the Goodstein
    price \citep{Gentzen1936}.
  \item[AOR cascade termination --- calibration, corpus tie.] The corpus's own cascade-termination
    theorem descends through the Dershowitz--Manna multiset extension of a well-founded status order;
    G2 promotes that internal usage to a first-class escrow currency
    \citep{DershowitzManna1979,Tsiokos2026AOR}.
\end{description}

The imported status of the proof-theoretic price and the bounded scope of the Goodstein lab are
consolidated in \Cref{sec:nonclaims}.

\subsection{\lawGThree}\label{sec:clusterA:g3}

Let $(S, \to)$ be a transition system and let a run prefix be a finite chain
$s_0 \to s_1 \to \cdots \to s_n$. Each transition carries a declared actual cost
$a_i = a(s_{i-1} \to s_i) \in \mathbb{Q}_{\ge 0}$, and a declared potential
$\Phi : S \to \mathbb{Q}_{\ge 0}$ induces the amortized cost
\[
  \widehat{a}_i \;=\; a_i + \Phi(s_i) - \Phi(s_{i-1}).
\]
G3 is a finite-prefix \emph{safety} law: it bounds total cost on every checked prefix and asserts
nothing about eventual success. The division of labor against G1 is exact --- G1 asks whether an
unbounded run eventually earns a positive descent certificate and may organize insolvent tails through
ghosts and completions; G3 assumes the potential is supplied and proves only what telescoping gives,
using no ghost, no completion, and no eventual-success clause. The Foundations~IV cousin F27
(Conservation as Orbit Descent) characterizes quantities that are orbit-constant; $\Phi$ is not
required to be orbit-constant and usually is not --- its lawful content is that potential
\emph{differences} telescope against visible costs. The row carries status tag \texttt{theorem}.

\begin{theorem}[\lawGThree]\label{thm:G3}
For every finite run prefix $s_0 \to \cdots \to s_n$,
\[
  \sum_{i=1}^{n} a_i
  \;=\; \sum_{i=1}^{n} \widehat{a}_i + \Phi(s_0) - \Phi(s_n)
  \;\le\; \sum_{i=1}^{n} \widehat{a}_i + \Phi(s_0),
\]
the inequality because $\Phi(s_n) \ge 0$. Consequently, if a declared class of runs has
$\widehat{a}_i \le B$ for every step, then every length-$n$ prefix satisfies
\[
  \sum_{i=1}^{n} a_i \;\le\; n B + \Phi(s_0).
\]
\end{theorem}

\begin{proof}
Expanding the amortized sum telescopes the potential:
\[
  \sum_{i=1}^{n} \widehat{a}_i
  \;=\; \sum_{i=1}^{n} \bigl(a_i + \Phi(s_i) - \Phi(s_{i-1})\bigr)
  \;=\; \sum_{i=1}^{n} a_i + \Phi(s_n) - \Phi(s_0).
\]
Rearranging gives the identity; nonnegativity of $\Phi(s_n)$ gives the bound; summing
$\widehat{a}_i \le B$ gives the consequence. That is the whole theorem --- every calibration instance
is an audit that the chosen $\Phi$ is nonnegative and the amortized costs uniformly
bounded.\footnote{Anchored to the substrate Lean at
\path|lean/SixBirdsFoundationsVI/Laws/G3AmortizedCurrency.lean|: \path|telescoping_identity|,
\path|telescoping_bound|, \path|amortized_sum_le_scale|, \path|uniform_actual_cost_bound|. The
mechanization realizes the algebraic theorem surface over exact integer-valued costs and potentials
with explicit nonnegativity hypotheses (\path|NonnegativePotential|, \path|NonnegativeActualCosts|);
the rational-valued statement above is the same telescoping algebra. Status tag \texttt{theorem}.}
\end{proof}

\paragraph{Currency legality.} A potential is a G3 \emph{currency} --- rather than a proxy --- only
when the declared audit contains all three of: $\Phi$ nonnegative on the state class; amortized costs
computed by the exact telescoping formula; and the resulting bound controlling a declared run-level
feasibility claim (total cost, or per-operation amortized cost). A non-telescoping score, or a
post-hoc statistic that merely correlates with cost, is a proxy and is rejected by the gate. This
clause is deliberately an audit discipline rather than a mechanized statement; the Lean file proves
the algebra, and legality is checked at the instance.

The audit records classify as follows (the table classifies a declared pair (run family, $\Phi$), not
the run family alone):

\begin{center}
\footnotesize
\begin{tabularx}{\textwidth}{@{}lXX@{}}
\toprule
Case & Status & Witness type \\
\midrule
(a) Prepaid & The audited $\Phi$ is legal and non-increasing at every checked step. &
Non-increasing-potential audit plus the telescoping bound. \\
(b) Borrow-and-repay & Some step increases $\Phi$, and $\widehat{a}_i$ stays uniformly bounded over
the family. & Credit ledger with bounded amortized costs; binary counters and dynamic arrays. \\
(c) Insolvent in the declared class & No legal $\Phi$ from the declared candidate class controls the
stated claim. & Lower-bound witness against the whole candidate class. \\
\bottomrule
\end{tabularx}
\end{center}

In interaction terms, G3 is P6 ledger content gated by P2 currency legality: the potential is the
balance, actual costs are visible debits, amortized costs are audited charges, and the gate rejects
proxies. The lab record: G3-L1 ran the binary-counter calibration (seed $7301$; $1000$ increments from
$44142$ to $45142$) and observed exact amortized cost $2$ at every step, and the dynamic-array
calibration (seed $7302$; $1000$ insertions after $3$ warmup insertions; $8$ resize steps) with
amortized cost at most $3$ everywhere and exactly $3$ at each resize --- both verifying the
telescoping identity exactly in integer arithmetic.

\subsection*{Layer instantiations}\label{layer-inst:G3}
\begin{description}
  \item[Binary-counter increment --- calibration.] $\Phi = $ number of $1$-bits gives exact amortized
    cost $2$ per increment \citep{Tarjan1985,CormenEtAl}.
  \item[Dynamic-array doubling --- calibration.] $\Phi = \max(0,\, 2 \cdot \mathrm{size} -
    \mathrm{capacity})$ gives amortized insertion cost at most $3$ under unit insert/copy costs
    \citep{CormenEtAl}.
  \item[Splay trees --- calibration.] The Sleator--Tarjan access lemma runs on the rank potential
    $\sum_x \log(\mathrm{size}(\mathrm{subtree}_x))$, with the constant depending only on the
    logarithm base; the instance is cited, not mechanized in the integer-valued Lean file
    \citep{SleatorTarjan1985}.
  \item[Passivity and storage functions --- structural.] Dissipativity theory bounds the change of a
    nonnegative storage function by supplied work; the corpus's existing passivity bridge-slot
    generalizes to the full finite-prefix currency law \citep{Willems1972,Tsiokos2026EmergenceCalculus}.
  \item[Energy-harvesting battery budgets --- structural.] Battery state is an explicit storage
    potential constraining consumption on every finite prefix under energy causality; the instance
    stays structural until one hardware model and cost convention are fixed
    \citep{OzelEtAl,YangUlukus}.
\end{description}

The absence of any potential-existence claim, the safety-not-liveness scope, and the integer-valued
mechanization surface are consolidated in \Cref{sec:nonclaims}.

% =============================================================================
\section{Cluster B --- \clusterB}\label{sec:clusterB}

Every obstruction in the earlier corpus is declared up front and holds still: a fixed hiddenness
interface, a fixed patch system, a fixed stability package. The four laws of this cluster cover the
obstructions that do not hold still --- covers that must move with the target (\lawGFour, G4),
prediction horizons that recede as the run advances (\lawGFive, G5), obstructions manufactured by the
run's own history (\lawGSix, G6), and obstructions selected adaptively by a budgeted adversary
(\lawGSeven, G7). One theme repeats: in each law, the positive certificate is stage-indexed or
strategy-valued data over an unbounded run, and the corresponding no-go is the demonstration that a
fixed, finite, or static package cannot substitute for it.

None of the four is a free theorem. G4 and G6 carry status tag \texttt{schema}; G5 and G7 carry
\texttt{calibration-anchored schema}, each anchored by a completely solved classical instance --- the
base-2 reverse-and-add orbit of $22$ for G5, and the angel problem at every power for G7.

\subsection{\lawGFour}\label{sec:clusterB:g4}

Let $T$ be an infinite index set and $\{P(t) : t \in T\}$ a target family. A \emph{certificate patch}
is data $c_\alpha$ validating $P(t)$ for every $t$ in a declared region $R_\alpha \subseteq T$; the
regions are part of the audit. A cover is \emph{moving} relative to a declared patch class when no
finite subfamily covers $T$ and the patch parameters required for targets in $T$ are unbounded. Two
equivalent audit shapes appear in instances: a global form, with uncovered remainders
$U_n = T \setminus \bigcup_{\alpha \le n} R_\alpha$ descending to empty under a declared exhaustion
order, and a pointwise form, with per-target residuals $\rho_0(t), \rho_1(t), \ldots$ and budgets
$b_j(t)$ descending strictly in a well-founded order until the residual empties.

The Foundations~IV cousin F4 (Local--Global Obstruction) classifies \emph{finite} patch systems:
finitely many local object sets, overlap compatibility, and a restriction map. G4 changes the axis
from finite gluing to unbounded exhaustion --- the patch family may grow with the target, the
certificate is stage-indexed, and closure is legal only when an explicit exhaustion ledger is
supplied. Local admissibility of patches never globalizes by itself. The row carries status tag
\texttt{schema}.

\begin{theorem}[\lawGFour]\label{thm:G4}
Declare a target family, a patch class, and a certificate language, and consider the hypotheses:
\begin{itemize}
  \item \textup{[H-G4-sound]} every patch certificate $c_\alpha$ proves $P(t)$ for every
    $t \in R_\alpha$ in the declared local certificate system;
  \item \textup{[H-G4-exhaust]} one of the declared audit forms proves total coverage: globally,
    every $t \in T$ leaves $U_n$ at some finite $n$; or pointwise, every $t \in T$ has empty residual
    $\rho_N(t)$ at some finite $N$;
  \item \textup{[H-G4-language-complete]} if $P(t)$ holds for all $t \in T$, the declared certificate
    language contains an exhaustion certificate witnessing that closure --- a meta-hypothesis about
    the language, never inferred from local compatibility;
  \item \textup{[H-G4-fixed-leak]} for a declared restricted class $\mathcal{C}$ of admissible patch
    regions, every finite subfamily from $\mathcal{C}$ misses at least one target.
\end{itemize}
Then:
\begin{enumerate}
  \item \emph{(Positive schema.)} \textup{[H-G4-sound]} and \textup{[H-G4-exhaust]} together give
    $P(t)$ for every $t \in T$.
  \item \emph{(Conditional biconditional.)} Relative to \textup{[H-G4-sound]} and
    \textup{[H-G4-language-complete]}, closure of $\{P(t)\}$ is \emph{equivalent} to the existence of
    an exhaustion certificate in the declared language.
  \item \emph{(Fixed-package no-go.)} Under \textup{[H-G4-fixed-leak]}, no finite package from
    $\mathcal{C}$ closes the claim --- a statement about that restricted class only, blocking nothing
    for a richer moving-cover family.
\end{enumerate}
\end{theorem}

\begin{proof}
The positive schema composes its two hypotheses: \textup{[H-G4-exhaust]} supplies, for each target,
the finite patch or terminal residual certificate, and \textup{[H-G4-sound]} converts it into
$P(t)$. The biconditional's forward direction is \textup{[H-G4-language-complete]} verbatim; its
reverse direction is the positive schema. The no-go is immediate from
\textup{[H-G4-fixed-leak]}.\footnote{Anchored to the substrate Lean at
\path|lean/SixBirdsFoundationsVI/Laws/G4MovingCoverExhaustion.lean|: \path|positive_schema|,
\path|conditional_biconditional|, \path|fixed_package_no_go|, over the setup predicates
\path|Sound|, \path|Exhausted|, \path|LanguageComplete|, \path|FixedLeak|. Status tag
\texttt{schema}.}
\end{proof}

The calibration is the Fibonacci--Sylvester greedy algorithm on proper rationals, where the whole
schema discharges with elementary arithmetic. For $p/q$ in lowest terms with $0 < p < q$, the greedy
step chooses $m = \lceil q/p \rceil$, subtracts $1/m$, and leaves remainder $(pm - q)/(qm)$. From
$m - 1 < q/p \le m$: the right inequality gives $pm - q \ge 0$ (no overshoot), and the left gives
$pm - q < p$, so the reduced numerator of the remainder is strictly smaller than $p$. The numerator
is therefore a strictly descending positive-integer budget --- the pointwise exhaustion form --- and
the expansion terminates in finitely many distinct unit fractions. The claim is deliberately scoped
to $(0,1)$ under the distinct-denominator convention. The fixed-package no-go is equally elementary:
a fixed finite denominator set has only finitely many subset sums, while $(0,1)$ contains infinitely
many rational targets, so no fixed finite package covers them; the greedy cover succeeds only because
its denominators move with the target.

The conjectural flagship is Erd\H{o}s--Straus: $4/n = 1/x + 1/y + 1/z$ for all $n \ge 2$. Prime
reduction is verified --- a solution for $4/p$ scales to $4/(mp)$, so a composite counterexample
would contain a prime one. Mordell-style modular identities are reported to cover all $n$ outside the
residue classes
\[
  n \;\equiv\; 1,\ 121,\ 169,\ 289,\ 361,\ 529 \pmod{840},
\]
and the quadratic-residue obstruction --- due to Mordell --- explains why no finite identity family
can finish the job: an identity for $n \equiv r \pmod{p}$ can exist only when $r$ is a
non-quadratic-residue modulo $p$, and $1$ is a quadratic residue modulo every $p > 1$, so the class
$n \equiv 1$ is never coverable by any finite identity package.\footnote{The mod-$840$ class list is
cross-verified against independent secondary sources; the specific polynomial identities themselves
are attributed to \citet[Ch.~30]{Mordell1969} but have not been independently retrieved from the
primary text, and are used here only as the fixed-package example, not as load-bearing input to any
theorem.} This is a concrete \textup{[H-G4-fixed-leak]} witness for the polynomial-identity class; it
neither proves nor refutes the conjecture. Density and average-count results
\citep{Vaughan1970,ElsholtzTao2013} remain \texttt{probabilistic} relative to the pointwise target.

The audit records classify by priority: (a) finitely coverable (a finite patch subfamily covers
$T$); (b) moving-cover-exhaustible (no finite package, but a global or pointwise exhaustion budget
--- the Fibonacci--Sylvester calibration); (c) fixed-package no-go for a declared restricted class
(coexisting, without contradiction, with a moving-cover proof in a richer class --- different audit
records); (d) probabilistically-covered-only (density-one, almost-all, or bounded-search evidence,
never upgraded to pointwise closure); (e) uncertified. In interaction terms, G4 is P2
feasibility-gating over an unbounded family plus a P6 exhaustion ledger; without the ledger, the
inference from local to global admissibility is blocked. The lab record: G4-L1 swept all $124750$
proper fractions $p/q$ with $0 < p < q \le 500$ ($76115$ distinct reduced targets); every greedy
expansion terminated with a strictly decreasing numerator budget and an exact unit-fraction
certificate --- the positive schema's calibration, exercised end to end.

\subsection*{Layer instantiations}\label{layer-inst:G4}
\begin{description}
  \item[Fibonacci--Sylvester Egyptian fractions --- calibration.] Target-dependent finite
    distinct-denominator certificates with strict numerator descent as the exhaustion budget
    \citep{Fibonacci,Eppstein1995}.
  \item[Erd\H{o}s--Straus --- conjectural flagship.] Open; prime reduction verified; the mod-$840$
    list and quadratic-residue obstruction are sub-results about identity families, carried with the
    primary-source flag above \citep{Mordell1969,ElsholtzTao2013}.
  \item[Erd\H{o}s covering systems --- calibration.] Finite congruence covering systems are the pure
    fixed-cover analogue; Hough's solution of the minimum-modulus problem bounds the least modulus in
    any distinct covering system \citep{Hough2015}.
  \item[Symbolic integration by pattern families --- structural.] Risch-style and rule-based
    integration grow expression-dependent case families rather than one finite pattern table; an
    instance requires the patch regions and exhaustion budget of a concrete integrator trace
    \citep{Risch1969,Bronstein2005}.
  \item[CEGAR and proof automation --- structural.] Counterexample-guided abstraction refinement
    exhausts spurious counterexamples with input-dependent precision; structural until a concrete
    verifier supplies the well-founded refinement budget \citep{ClarkeEtAl2003}.
\end{description}

The Erd\H{o}s--Straus nonclaims and the probabilistic status of density evidence are consolidated in
\Cref{sec:nonclaims}.

\subsection{\lawGFive}\label{sec:clusterB:g5}

Fix a base $b \ge 2$. For a positive integer $n$ with canonical digit string, let $\mathrm{rev}_b(n)$
be the integer represented by the reversed string and $R_b(n) = n + \mathrm{rev}_b(n)$ the
reverse-and-add map; the target predicate is usually $\mathrm{Pal}_b(n)$, palindromicity. The
depth-$m$ current quotient is $q_m(n) = n \bmod b^m$ --- the last $m$ digits. For a declared future
probe $F_j$ (such as $q_m(R_b^{\,j}(n))$, or the palindrome verdict within $j$ steps), the
\emph{carry horizon} $h_j(x)$ is the least depth $M$ such that all $y$ with $q_M(y) = q_M(x)$ give
the same probe value as $x$, and $\infty$ when no finite depth suffices. The law quantifies over an
unbounded run: can a fixed quotient depth remain predictive as $j \to \infty$? A \emph{pattern
certificate} is a regular or automatic language $\mathcal{P}$ over digit strings, with declared phase
convention, transition closure, and target exclusion.

The Foundations~IV cousins are again single-instant: F13a fixes one hiddenness interface and one
exposure test; F3 compares two routes once. G5 is their step-indexed strengthening --- a
\emph{sequence} of quotient-prediction failures as the horizon grows along the orbit, against a
confinement certificate that must propagate through the entire run. The row carries status tag
\texttt{calibration-anchored schema}.

\begin{theorem}[\lawGFive]\label{thm:G5}
\textbf{No-go schema.} Assume
\begin{itemize}
  \item \textup{[H-G5-unbounded-horizon]} for every fixed depth $m$ there are a step $j$ and states
    $y, z$ in the declared carrier with $q_m(y) = q_m(z)$ but different declared future probe values
    --- a nonempty step-indexed split-pair obstruction $\Delta(q_m, R_b^{\,j}, r_j) \neq \varnothing$.
\end{itemize}
Then no fixed-depth quotient $q_m$ is future-sufficient for the declared unbounded run: for every
$m$, the future probe fails to factor through $q_m$.

\textbf{Confinement schema.} Assume
\begin{itemize}
  \item \textup{[H-G5-closed-pattern]} a declared regular or automatic language $\mathcal{P}$
    satisfies $R_b(\mathcal{P}) \subseteq \mathcal{P}$ (cyclically $R_b(P_i) \subseteq P_{i+1}$ when
    phase-indexed, with $\mathcal{P} = \bigcup_i P_i$), and no member of $\mathcal{P}$ satisfies the
    target predicate.
\end{itemize}
Then if some iterate $R_b^{\,N}(x)$ lies in $\mathcal{P}$, no later iterate satisfies the target:
closure gives $R_b^{\,N+t}(x) \in \mathcal{P}$ for every $t \ge 0$, and target-freeness gives
$\lnot\,\mathrm{Pal}_b(R_b^{\,N+t}(x))$ for every $t \ge 0$.

Both halves are conditional: the no-go does not by itself prove escape from the target, and the
confinement applies only after a concrete target-free closed language has been verified.
\end{theorem}

\begin{proof}
The no-go half is the split-pair definition applied at each depth: in reverse-and-add, the low $m$
digits after a step depend on high digits brought down by the reversal, and over an unbounded orbit,
longer strings expose carries from outside every fixed suffix window ---
\textup{[H-G5-unbounded-horizon]} packages that moving-carry fact as an auditable obstruction family
rather than a heuristic. The confinement half is induction along the run: entry into $\mathcal{P}$
plus closure keeps every later iterate inside, and target-freeness excludes the
target.\footnote{Anchored to the substrate Lean at
\path|lean/SixBirdsFoundationsVI/Laws/G5CarryHorizonConfinement.lean|:
\path|no_fixed_depth_future_sufficient| and \path|no_future_sufficient_depth| over
\path|UnboundedHorizon|/\path|FutureSufficient|; \path|closed_pattern_persists| and
\path|closed_pattern_confinement| over \path|ClosedUnderR|/\path|TargetFree|/\path|ClosedPattern|.
Status tag \texttt{calibration-anchored schema}.}
\end{proof}

The calibration is the solved base-2 instance: the orbit of $22 = 10110_2$ never reaches a binary
palindrome. Directly, $10110 \to 100011 \to 1010100 \to 1101001 \to 10110100$, and from there the
orbit enters a four-phase regular family, defined for $r \ge 2$ by
\[
  P_0(r) = 10\,1^r\,01\,0^r, \quad
  P_1(r) = 11\,0^{r-2}\,1\,000\,1^{r-2}\,01, \quad
  P_2(r) = 10\,1^r\,01\,0^{r+1}, \quad
  P_3(r) = 11\,0^r\,10\,1^{r-1}\,01,
\]
with the phase cycle, verified by binary addition with carries,
\[
  R(P_0(r)) = P_1(r), \qquad
  R(P_1(r)) = P_2(r), \qquad
  R(P_2(r)) = P_3(r), \qquad
  R(P_3(r)) = P_0(r+1).
\]
No string in the family is a palindrome: $P_0$ and $P_2$ begin with $1$ and end with $0$; $P_1$ and
$P_3$ begin with $11$ and end with $01$, so second and penultimate digits differ. The four pre-entry
strings are also non-palindromic, so the whole orbit is certified target-free --- discharging
\textup{[H-G5-closed-pattern]} with the full four-phase family as the invariant (note the closure is
cyclic phase-to-phase; the frequently quoted pattern $10\,1^k\,01\,0^k$ recurs under $R^4$, every
fourth iterate, not under single steps).\footnote{The base-2 result is recorded at OEIS A060382 and
attributed to Kevin Brown's MathPages note; no peer-reviewed publication of the proof was located, so
the attribution is informally sourced --- and this Kevin Brown is distinct from the Cornell
topologist Kenneth~S.~Brown, with whom secondary discussions sometimes conflate him
\citep{OEIS-A060382,Brown-MathPages}.} For base 10, no analogous certificate is known: $196$ remains
open, and finite searches, however deep, are bounded evidence only \citep{Walker-Fourmilab}.

The audit records classify by priority: (a) confined-to-target-basin (explicit hitting time or basin
certificate); (b) certified escape (a target-free closed pattern family --- the base-2 $22$
calibration); (c) horizon-budgeted-undetermined (bounded evidence only --- base-10 $196$ lives
here). In interaction terms, G5 combines a P1 descent obstruction with P5 packaging: the split-pair
family records that fixed quotients fail as predictive carriers, and the regular language packages an
infinite target-free residue into a finite certificate; carry-horizon growth is hidden predictive
content, not entropy production. The lab record: G5-L1 followed the seed $22$ through the exact
pre-entry segment and then $20$ full phase cycles --- $80$ reverse-and-add transitions across $81$
phase states, $84$ steps total, largest string $48$ bits --- verifying every phase identity in exact
integer arithmetic. The lab calibrates the concrete instance of the hypothesis over a bounded run;
the universal confinement statement is the Lean theorem, which takes the closed-pattern hypothesis as
supplied data.

\subsection*{Layer instantiations}\label{layer-inst:G5}
\begin{description}
  \item[Base-2 reverse-and-add --- calibration.] The seed $22$ is certified palindrome-free by the
    four-phase invariant above \citep{OEIS-A060382,Brown-MathPages}.
  \item[Base-10 Lychrel candidates --- conjectural.] The $196$ orbit has extensive finite computation
    and no proof either way \citep{Walker-Fourmilab}.
  \item[Carry-propagation side channels --- structural.] Variable carry propagation can expose
    digit-position information through timing or power; a G5-style moving-horizon risk once the
    hardware model declares readout and leakage predicates \citep{Kocher1996}.
  \item[Finite-precision error fronts --- structural.] Rounding and cancellation move the depth at
    which hidden low-order information affects later results; structural until a specific algorithm
    supplies the quotient chain and horizon audit \citep{Higham2002}.
  \item[Automatic sequences and regular invariants --- calibration/structural.] Regular languages and
    finite automata are the certificate class of the base-2 proof; other digit systems join only when
    transition closure and target exclusion are checked \citep{AlloucheShallit2003}.
\end{description}

The Lychrel nonclaims and the bounded-evidence discipline for deep base-10 searches are consolidated
in \Cref{sec:nonclaims}.

\subsection{\lawGSix}\label{sec:clusterB:g6}

Let $X$ be a countable carrier and let a run be $x_{n+1} = T(x_n, S_n)$ with
$S_n = G(x_0, \ldots, x_n) \subseteq X$: the \emph{endogenous obstruction set}, computed from the
orbit's own past by a declared history functional --- not fixed before the run. For a finite audit
window $W \subseteq X$, write $V_n = \{x_0, \ldots, x_n\}$ for the visited set and
$H_n(W) = W \setminus V_n$ for the unvisited holes. With a declared finite measure $\mu$ on windows,
the rate quantities are the obstruction growth $\gamma_n(W) = \mu((S_{n+1} \setminus S_n) \cap W)$
and a declared legal first-visit pressure $\rho_n(W)$ from $(x_n, S_n)$ toward $H_n(W)$; the exact
form of $\rho_n$ is instance data (a $0/1$ legal-return indicator, a probability, a channel count),
but the audit is not optional --- $\rho_n$ is computed before any future visit is used as evidence.

A terminological note is required here, because the corpus contains a near-collision. In this law,
\emph{endogenous} means that the obstruction is manufactured by the orbit's own history and analyzed
externally, by the theorist auditing the run. It is distinct from the sense of ``endogenous'' in
Foundations~V's ``Endogenous Closure,'' where a system-side apparatus performs its own closure and
repair \citep{FoundationsVCognition}. G6 never credits the system with certifying itself: all G6
certificates are theorist-side audits of the generated obstruction landscape, and that boundary
(NC-16/NC-17 in the corpus's numbering) is load-bearing. The Foundations~IV cousin F6 (No-Needles
Stability) declares its stability package --- currencies, budgets, transport selector --- up front
and kills unpriced needles against it; G6 reverses the provenance: the needle is manufactured by the
run, and the theorist audits the race between return pressure and growing obstruction after the
fact. The row carries status tag \texttt{schema}.

\begin{theorem}[\lawGSix]\label{thm:G6}
\textbf{Covering schema.} Assume
\begin{itemize}
  \item \textup{[H-G6-dominant-pressure]} there are a finite exception set $E \subseteq X$, a cofinal
    exhaustion $W_1 \subseteq W_2 \subseteq \cdots$ with $\bigcup_M W_M = X \setminus E$, and for
    each $M$ a finite block decomposition of times after some $N_M$ such that whenever $H_n(W_M)$ is
    nonempty at the start of a block $I$, the audited inequality
    \[
      \sum_{i \in I} \rho_i(W_M) \;\ge\; 1 + \sum_{i \in I} \gamma_i(W_M)
    \]
    holds and the excess is charged to at least one distinct first visit in $H_n(W_M)$ during $I$.
\end{itemize}
Then every point of $X \setminus E$ is eventually visited: the orbit is cofinite-covering relative to
$E$.

\textbf{Hole-forming schema.} Assume
\begin{itemize}
  \item \textup{[H-G6-dominant-obstruction]} there are a target $y \in X$, a time $N$ with
    $y \notin V_N$, and a declared sequence of separating cuts $C_n(y)$ such that for every
    $n \ge N$ every legal path from $(x_n, S_n)$ to $y$ crosses $S_n$ or a certified blocked cut,
    with the quantitative excess $\sum_{i \in I} \gamma_i(C_i(y)) > \sum_{i \in I} \rho_i(C_i(y))$ on
    every later audit block, tied to the displayed cut rather than inferred from density.
\end{itemize}
Then $y$ is a permanent hole: $y \notin V_n$ for all $n \ge N$.

Both directions are schemas: no concrete system satisfies either hypothesis until its rate ledger and
witness cuts are supplied, and the two hypotheses can only conflict on a single target if the audit
record itself is inconsistent.
\end{theorem}

\begin{proof}
For covering, fix a window $W_M$; it is finite. Each audited block that starts with a nonempty hole
set buys at least one distinct first visit in that hole set, so $|H_n(W_M)|$ strictly decreases over
charged blocks and cannot decrease forever; hence $W_M$ is covered after finitely many blocks, and
cofinality of the exhaustion covers all of $X \setminus E$. For hole-forming, every legal route to
$y$ after $N$ crosses an obstruction whose audited growth dominates the available pressure on every
block, so no certified legal visit to $y$ can occur after $N$; since $y \notin V_N$, it never
appears.\footnote{Anchored to the substrate Lean at
\path|lean/SixBirdsFoundationsVI/Laws/G6EndogenousNeedleGeneration.lean|: \path|covering_schema| over
\path|DominantPressure|, and \path|hole_forming_schema| over \path|DominantObstruction|. Status tag
\texttt{schema}.}
\end{proof}

The conjectural flagship is Recam\'an's sequence: $a_0 = 0$ and $a_n = a_{n-1} - n$ when that value
is positive and unused, else $a_n = a_{n-1} + n$; the visited set is the orbit-generated obstruction
set. Whether every positive integer eventually appears is open; the strongest computational fact is
that $852655$ remains the smallest missing value after computations extending to $10^{612}$ terms
--- a bounded census, not a permanent-hole proof \citep{OEIS-A005132,OEIS-A057167}. The status
justification is recorded honestly: the two schemas are each calibrated by a solved \emph{designed}
variant (below), but by deliberately tuned toys rather than by a recognized classical instance, and
the grade therefore stays \texttt{schema} rather than \texttt{calibration-anchored} --- the contrast
with G7 and G9, which anchor on the angel problem and rule 184, is intentional.

The classification runs on two axes. System-level: (S-a) cofinite-covering relative to $E$
(\textup{[H-G6-dominant-pressure]} discharged on a cofinal exhaustion); (S-b) uncontrolled
obstruction growth (the audit shows growth escaping the pressure ledger, without a covering
certificate); (S-c) system-uncertified --- where Recam\'an currently sits. Per-target: (T-a) covered
(hitting time, or covered by a system certificate with $y \notin E$); (T-b) certified permanent hole
(\textup{[H-G6-dominant-obstruction]} discharged for this $y$); (T-c) target-uncertified. In
interaction terms, G6 is P5 packaging of the history-generated obstruction plus P2 feasibility gating
through the legal-move predicate, with a P6 rate ledger for $\gamma_n$ and $\rho_n$ --- always
theorist-side. The lab record spans both axes. G6-L1 ran a $10^7$-step exact census of Recam\'an:
bounded-window coverage high but incomplete, return-pressure proxy exactly $4999986/10^7$ (the
fraction of steps where the backward move was legal), and a visible smallest-missing trajectory ---
an (S-c) exhibit, instantiating no schema. G6-L2 supplies the two solved designed variants: variant~A
(jump $2n$) visits only even values --- $10^6$ steps, all-even invariant proved by parity induction,
$736749$ distinct values visited, zero odd visits --- so every odd positive integer is a certified
permanent hole, discharging the hole-forming schema's shape; variant~B (constant jump $1$) satisfies
$a_n = n$ with visited set exactly $\{0, \ldots, n\}$, proved by induction, discharging the covering
schema's shape.

\subsection*{Layer instantiations}\label{layer-inst:G6}
\begin{description}
  \item[Recam\'an's sequence --- conjectural flagship.] Coverage open; $852655$ long-missing in
    bounded censuses \citep{OEIS-A005132,OEIS-A057167}.
  \item[Self-avoiding walks --- structural.] The walk's past forbids future sites --- an endogenous
    obstruction landscape; structural because the coverage/hole dichotomy needs a rate ledger beyond
    the standard model \citep{MadrasSlade1993,Flory1953}.
  \item[Novelty search and curiosity-driven exploration --- structural.] Visited-state penalties and
    pseudo-count bonuses let history reshape the exploration landscape. The seam with Foundations~V
    is explicit here: V prices the agent's own probe-acquisition economy from the system side, while
    G6 audits the orbit's obstruction set from the theorist's side
    \citep{BellemareEtAl2016,FoundationsVCognition}.
  \item[Online algorithms with irrevocable commitments --- structural.] Accepted commitments alter
    the feasible set for later requests; structural until a scheduler supplies $\gamma_n$, $\rho_n$,
    and a covering or hole witness \citep{BorodinElYaniv1998,BoyarEtAl}.
  \item[Greedy sequences in combinatorial number theory --- structural.] Mian--Chowla and
    Stanley-type sequences avoid constraints manufactured by their own earlier terms; endogenous
    generation without the rate audit \citep{MianChowla1944,Rolnick}.
\end{description}

The Recam\'an nonclaims, the Phase-1 status of the two rate hypotheses, and the self-certification
boundary are consolidated in \Cref{sec:nonclaims}.

\subsection{\lawGSeven}\label{sec:clusterB:g7}

Let $B$ be a homogeneous board --- in the standard game, $\mathbb{Z}^2$ with the Chebyshev metric
$d_\infty$. Fix an agent mobility parameter $p \ge 1$ and a devil deletion budget $b \ge 1$. In the
angel game, a power-$p$ angel moves each turn to a different undeleted square within
$d_\infty \le p$ (jumping over deleted squares if necessary); the devil then deletes at most $b$
visible squares not currently occupied. With $D_n$ the deleted set and $a_n$ the angel position, the
visible finite history is $H_n = (a_0, D_0, \ldots, a_n, D_n)$. The angel wins if it can move
forever; the devil wins if it can force a finite turn with no legal angel move. The adversary's
strategy space is declared game data --- no hidden-oracle devil is ever implicitly assumed.

The Foundations~IV cousin F6 declares its package up front and has no turn index and no reacting
adversary; G7 adds the game-valued axis the corpus lacks: obstruction selected adaptively by a
visible adversary under a per-turn budget, certified by strategy proofs over unbounded play. The row
carries status tag \texttt{calibration-anchored schema}.

\begin{theorem}[\lawGSeven]\label{thm:G7}
Fix a board, mobility rule, and devil budget, and consider:
\begin{itemize}
  \item \textup{[H-G7-escape-strategy]} a declared angel strategy $\sigma_A$ from finite visible
    histories to legal moves, with an invariant $I_n(H_n)$ and a proof that against every legal devil
    strategy, $\sigma_A$ preserves $I_n$ and leaves a legal undeleted move at every finite turn;
  \item \textup{[H-G7-confinement-strategy]} a declared devil strategy $\sigma_D$ from finite visible
    histories to at most $b$ legal deletions, with a wall, cut, ranking, or amortization proof that
    against every legal angel strategy, $\sigma_D$ forces a finite turn $N$ with no legal angel move.
\end{itemize}
Exactly these certificate statuses are admissible:
\begin{enumerate}
  \item if \textup{[H-G7-escape-strategy]} is discharged, the angel wins --- every play consistent
    with $\sigma_A$ is infinite;
  \item if \textup{[H-G7-confinement-strategy]} is discharged, the devil wins --- every play against
    $\sigma_D$ is trapped in finitely many turns;
  \item if neither is supplied, the parameter lies in the declared ruleset's open band.
\end{enumerate}
The two certificates are mutually exclusive for a fixed declared game: playing the two strategies
against each other would produce a single play both infinite and finitely trapped.

For the classical game on $\mathbb{Z}^2$ with $b = 1$ the threshold is fully known: the power-$1$
angel loses, and every power-$p$ angel with $p \ge 2$ wins; no open band remains for the standard
two-dimensional one-devil game.
\end{theorem}

\begin{proof}
G7 does not reproduce the classical proofs; it extracts and types their certificates. On the
confinement side, the devil certificate is a visible wall-building schedule: in the power-$1$ case
the angel is too slow to cross a suitably amortized enclosing construction, and the proof obligation
is a finite strategy plus an invariant confining every angel response to a closable legal region ---
the early Conway--Berlekamp result \citep{BerlekampConwayGuy1982,Conway1996}. On the escape side, the
certificate is a route-maintenance invariant whose progress outpaces the unit deletion budget:
Kloster's constructive path-management strategy for the $2$-angel charges path updates to newly
blocked squares \citep{Kloster2007}; M\'ath\'e's proof reduces the real devil to a ``nice devil'' and
runs a wall-following maze invariant \citep{Mathe2007}; Bowditch's independent proof wins with a
$4$-angel through auxiliary game variants \citep{Bowditch2007}; G\'acs gives a hierarchical
high-power escape \citep{Gacs2007}. By monotonicity, a higher-power angel simulates any lower-power
winning strategy. These discharge the two hypotheses at the stated powers; the mutual-exclusivity
clause is the elementary play-off argument.\footnote{Anchored to the substrate Lean at
\path|lean/SixBirdsFoundationsVI/Laws/G7AdversarialMobilityConfinement.lean|:
\path|escape_never_stuck| over \path|EscapeStrategy|; \path|confinement_forces_trap| (via
\path|confinement_forces_trap_with_fuel|) over \path|ConfinementStrategy| and \path|AngelHonest|;
\path|escape_confinement_mutually_exclusive|. The \path|ConfinementStrategy| structure carries an
explicit devil-legality field, so the confinement certificate quantifies only over legal devil
schedules. Status tag \texttt{calibration-anchored schema}.}
\end{proof}

The game records classify as: (a) confined (\textup{[H-G7-confinement-strategy]} discharged ---
standard $p = 1$, $b = 1$); (b) escaping (\textup{[H-G7-escape-strategy]} discharged --- standard
$p \ge 2$, $b = 1$); (c) open band (neither supplied --- none remains for the standard game, but
variant boards and budgets may have open bands). In interaction terms, G7 is P2 feasibility/mobility
under a dynamic legal-move predicate plus P6 budget accounting of the deletion rate; the adversarial
structure is visible and typed in the audit. The lab record is deliberately modest: G7-L1 ran exact
bounded-horizon minimax on small finite tori ($8$ parameter settings). The one unconfounded
comparison, $N = 4$, $b = 2$, showed the expected qualitative threshold --- $p = 1$ trapped within
the horizon, $p = 2$ surviving to it --- while the $3 \times 3$ rows are geometric degeneracies
($p = 1$ and $p = 2$ have identical legal moves there) and $N = 4$, $b = 1$ was inconclusive at
horizon $8$. A finite torus is not the infinite board: survival to horizon means forced survival for
the tested turns only, and the exhibit is a qualified finite-board calibration, not evidence about
infinite play.

\subsection*{Layer instantiations}\label{layer-inst:G7}
\begin{description}
  \item[Angel problem --- calibration.] On $\mathbb{Z}^2$ with one deletion per turn: $p = 1$
    confined, $p \ge 2$ escapes
    \citep{BerlekampConwayGuy1982,Conway1996,Bowditch2007,Mathe2007,Kloster2007,Gacs2007}.
  \item[Pursuit-evasion and cop number --- structural/calibration.] Cops-and-robbers games classify
    mobility against capture strategies on declared graphs, with thresholds depending on graph class
    and capture rules \citep{NowakowskiWinkler1983,AignerFromme1984}.
  \item[Routing under adversarial link failure --- structural.] Adversarial queueing asks whether
    rerouting outpaces adversarial blockage; structural until a concrete network model supplies both
    strategies \citep{BorodinEtAl2001}.
  \item[Percolation navigation above criticality --- structural.] Supercritical percolation provides
    infinite open clusters through random obstruction fields --- a probabilistic analogue, not an
    adaptive devil, until the adversary model is declared \citep{Grimmett1999,Kesten1982}.
  \item[Robust motion planning with dynamic obstacles --- structural.] Explicit dynamic-obstacle
    models with safety and viability certificates become G7-like when the obstacle budget and planner
    strategy are declared \citep{LaValle2006,FraichardAsama2004}.
\end{description}

The board-relativity of thresholds, the visible-adversary requirement, and the exhibit-only status of
finite-board searches are consolidated in \Cref{sec:nonclaims}.
% =============================================================================
\section{Cluster C --- \clusterC}\label{sec:clusterC}

The corpus handles route defects one comparison at a time: two routes disagree, the disagreement is
recorded as a binary witness, and repair either succeeds or fails at that single square. Nothing in the
earlier catalogs integrates route data into a potential over an unbounded run, derives motion from
substrate dynamics, or asks which declared content can survive iterated loss. The three laws of this
cluster occupy exactly that verified gap. \lawGEight{} (G8) identifies the hypothesis under which the
firing counters of \emph{every} legal route through an unbounded run integrate into one conserved
odometer, and the further hypothesis under which that odometer is a least-action ledger.
\lawGNine{} (G9) supplies the census discipline under which persistent local defects evacuate into
certified transport records --- motion emerging from, rather than being declared over, the substrate.
\lawGTen{} (G10) proves that one declared boundary readout can persist through every iterate of a
lossful map, but only because a separately declared interior invariant regenerates the protection at
each step.

The grades are deliberately mixed. G8 and G10 carry status tag \texttt{theorem}. G9 carries
\texttt{calibration-anchored schema}: its census and transporter records are typed and mechanized, and
the rule-184 traffic automaton discharges the entire mechanism as solved mathematics, but the abstract
implication still passes through named census and absorption hypotheses rather than a carrier-agnostic
proof.

\subsection{\lawGEight}\label{sec:clusterC:g8}

Let $C$ be a configuration space and $V$ a finite or locally finite index set of sites. For each
$v \in V$, let $m_v : C \to C$ be a partial local move guarded by a legality predicate $L_v(c)$. A run
from $x \in C$ is an unbounded legal sequence
\[
  x = c_0 \xrightarrow{\,v_0\,} c_1 \xrightarrow{\,v_1\,} c_2 \xrightarrow{\,v_2\,} \cdots
\]
continued until no legal move remains, if such a time exists. Its \emph{counter vector} is
\[
  u_{\mathrm{run}}(v) \;=\; \#\{\, i : v_i = v \,\} \;\in\; \mathbb{N}^{V},
\]
finite in every coordinate when the run terminates. A configuration is \emph{stable} when no $L_v$
holds, and a vector $n \in \mathbb{N}^{V}$ is \emph{sufficient} for $x$ when the multiset of moves it
prescribes reaches a stable configuration in some legal admissible stabilizing order.

The temporal quantifier is the content of the law. G8 quantifies over every legal firing order of an
unbounded run and over the whole accumulated counter vector --- not over one fixed comparison of two
routes. Its closest Foundations~IV cousins mark the contrast. F3 (Holonomy--Memory Repair) compares
exactly two fixed routes once and, when they disagree at the predictive quotient, records a residue;
it never produces a global scalar or a canonical counter ledger. F27 (Conservation as Orbit Descent)
presumes that a conserved quantity already exists and characterizes conservation as orbit-constancy;
it never constructs the quantity. G8 does what neither row can: under an abelian compatibility
hypothesis and termination, it \emph{constructs} the conserved ledger from legal move counters and
proves its route independence. The row carries status tag \texttt{theorem}; both parts below are the
carrier-agnostic forms of classical abelian-sandpile results, with proofs adapted from
\citet{Dhar1990} and \citet{FeyLevinePeres2010}.

\begin{theorem}[\lawGEight]\label{thm:G8}
Let $(C, V, \{m_v\}, L)$ be a move system with initial configuration $x$, and assume:
\begin{itemize}
  \item \textup{[H-G8]} (abelian compatibility) for every configuration $c$ and all \emph{distinct}
    sites $v \neq w$: if $L_v(c)$ and $L_w(c)$ both hold, then firing either move preserves legality
    of the other, the two composites agree,
    \[
      m_w(m_v(c)) \;=\; m_v(m_w(c)),
    \]
    and both two-step routes record the same counter increment $e_v + e_w$. The case $v = w$ is
    outside the scope of this hypothesis: \textup{[H-G8]} does not require a move to preserve its own
    legality after firing.
  \item \textup{[termination at $x$]} every maximal legal run from $x$ terminates.
\end{itemize}
Then:
\begin{enumerate}
  \item \emph{Order-independent stabilization.} All legal terminating runs from $x$ end at the same
    stable configuration $x^{\circ}$.
  \item \emph{Odometer invariance.} All legal terminating runs from $x$ have the same counter vector
    $u_x \in \mathbb{N}^{V}$.
\end{enumerate}
If, in addition,
\begin{itemize}
  \item \textup{[H-G8-mono]} (least-action monotonicity) the move system carries a declared comparison
    structure under which stabilizing scripts are compared pointwise, legal move maps are compatible
    with the comparison, and adding available resource or delaying commuting legal moves cannot make a
    genuinely necessary firing disappear from every stabilizing script,
\end{itemize}
then the odometer is pointwise minimal among legally sufficient scripts: for every $n \in
\mathbb{N}^{V}$ whose prescribed multiset of moves is realized by some legal admissible stabilizing
order from $x$,
\[
  u_x(v) \;\le\; n(v) \qquad \text{for every } v \in V.
\]
This is a legal-versus-legal least-action statement. It does not compare the odometer against the
fully general Fey--Levine--Peres class of arbitrary nonnegative stabilizing scripts, which includes
force-firing scripts passing through illegal intermediate configurations.
\end{theorem}

\begin{proof}
The abelian part is the sandpile diamond argument abstracted away from graphs. \textup{[H-G8]} is a
local commutation rule with matched counters: whenever two distinct legal moves are simultaneously
available, either order reaches the same configuration and records the same two increments. Any two
legal terminating orders can therefore be compared by repeatedly swapping adjacent independent legal
moves, without changing the multiset of moves already accounted for; the termination hypothesis
excludes the infinite-postponement defect that would otherwise let a needed firing be delayed forever.
Local commutation thus propagates to global agreement of both the final stable configuration and the
firing counts.

Least action uses \textup{[H-G8-mono]}, not commutation alone. Compare a legal stabilization with a
legally sufficient vector $n$, and suppose the legal run fires some site more often than $n$
prescribes; take the first prefix at which the pointwise domination fails. At that prefix, abelian
compatibility together with the declared monotone comparison structure lets the already-fired legal
multiset be commuted against the sufficient multiset, contradicting the assumption that $n$ already
supplied enough moves to stabilize. The comparison never leaves the class of legally realizable
scripts; extending it to arbitrary force-fired scripts would require unconditional, rather than
legality-guarded, commutativity, which is not assumed.\footnote{Anchored to the substrate Lean at
\path|lean/SixBirdsFoundationsVI/Laws/G8OdometerAbelianization.lean|: Part~A is
\path|partA_order_independent_and_odometer| (with \path|order_independent_stabilization| and
\path|odometer_invariance|), Part~B is \path|least_action|. Status tag \texttt{theorem}.}
\end{proof}

The case enumeration separates the law from the weaker rewriting properties it is most easily confused
with:

\begin{center}
\footnotesize
\begin{tabularx}{\textwidth}{@{}lXX@{}}
\toprule
Case & Status & Witness type \\
\midrule
(a) Abelian with odometer & \textup{[H-G8]} holds and every maximal legal run from $x$ terminates;
with \textup{[H-G8-mono]}, least action is included. & Abelian-compatibility proof, termination
certificate, stable state $x^{\circ}$, odometer $u_x$, and (when claimed) the monotone comparison
proof. \\
(b) Confluent without canonical counters & The system terminates with a unique normal form, but
counter vectors depend on the route. & The explicit witness below, mechanized as
\path|caseB_not_abelian| and \path|caseB_route_dependent_counter|. \\
(c) Non-terminating & Some legal run from $x$ is infinite; no finite odometer exists to certify. &
Infinite legal run or failed termination budget. A nonterminating run with a separate finite currency
belongs to G3's amortized territory (\Cref{sec:clusterA:g3}). \\
(d) Terminating, non-confluent & All runs terminate but reach distinct stable configurations. & The
two-branch witness mechanized as \path|caseD_nonconfluent_witness|. \\
\bottomrule
\end{tabularx}
\end{center}

Case (b) is the separating witness that distinguishes odometer invariance from mere confluence, and it
is small enough to state in full. Take states $S, A, B, N$ with $N$ normal and legal moves
\[
  r : S \to A, \qquad s : S \to B, \qquad p : A \to N, \qquad q : B \to N, \qquad t : A \to N,
\]
and no others. Every move strictly decreases the rank $S > A, B > N$, so the system terminates, and
every maximal route from $S$ ends at the unique normal form $N$. But the three maximal routes carry
three different counter vectors: $e_r + e_p$, $e_r + e_t$, and $e_s + e_q$. Unique normal form does
not produce a canonical move-count ledger. The missing property is exactly \textup{[H-G8]}: after $r$
and $s$ compete at $S$, the unchosen branch is not preserved as a legal commuting move. Confluence
records that route defects can be \emph{joined}; G8 requires the strictly stronger compatibility that
\emph{integrates} route data into one odometer.

In interaction terms, G8 turns P3 route data into P6 audit content through a P2 conservation gate: the
legal firing order is route data, the counter vector is bookkeeping, and \textup{[H-G8]} is the gate
that makes the bookkeeping conserved across all legal routes. The claim is conditional on that gate
and on termination; it is not an unconditional inference from route mismatch to potential. The
computational record exercises both sides. Lab G8-L1 ran $500$ seeded random stabilizations on square
sandpile grids of sizes $10, 20, 30, 40, 50$ ($100$ random legal orders per configuration) and found
every run agreeing with the deterministic baselines in both final configuration and odometer vector.
Lab G8-L2 sampled $1000$ random legal orders of the case~(b) witness: every run reached $N$, and
exactly three distinct counter vectors appeared --- confluence without a canonical ledger, as the law
predicts.

\subsection*{Layer instantiations}\label{layer-inst:G8}
\begin{description}
  \item[Abelian sandpiles and rotor-routing --- calibration.] Finite graph sandpiles with a sink
    discharge \textup{[H-G8]} by Dhar's abelian property and the least-action part by the
    Fey--Levine--Peres monotonicity argument \citep{Dhar1990,FeyLevinePeres2010}. Rotor-routing is the
    deterministic sibling in the same abelian-network family \citep{HolroydEtAl2008}.
  \item[Quantitative confluence certificates in rewriting --- structural.] The corpus rewriting paper
    leaves the full confluence-as-flatness theorem package for finite terminating left-linear systems
    as deferred packaging; G8 supplies part of the missing quantitative layer, but only for abelian
    sub-cases where route counters commute --- it is not a general solution to that deferred package
    \citep{Tsiokos2026Rewriting}.
  \item[Distributed load balancing and gossip averaging --- structural.] Diffusive load balancing and
    randomized gossip conserve total mass or average under order-sensitive local exchanges; G8 applies
    to the subfamilies where the exchange counters commute and a finite stabilization surrogate exists
    \citep{Cybenko1989,BoydEtAl2006}.
  \item[Neural avalanche models --- structural, unverified for G8.] Neuronal avalanches supply a real
    cascade literature \citep{BeggsPlenz2003}, but biological avalanche dynamics are not automatically
    abelian and their event counters need not be route-independent; upgrading this row beyond
    structural requires a precise model with commuting local relaxations and a certified finite
    odometer.
  \item[Settlement netting and clearing vectors --- structural.] Clearing systems have ledger-like
    conservation and least-clearing fixed points, but ordinary settlement is not automatically an
    abelian local-move system; a strict instance must exhibit commuting settlement moves whose clearing
    vector is the odometer-like least-action object \citep{EisenbergNoe2001}.
\end{description}

The classical-novelty caveat, the open least-action computational exhibit, and the scope of the
legal-versus-legal comparison are consolidated in \Cref{sec:nonclaims}.

\subsection{\lawGNine}\label{sec:clusterC:g9}

Let $X$ be a spatial lattice or graph carrying a translation action $\mathrm{shift}_d$, let $S$ be a
finite local state set, and let $T : S^{X} \to S^{X}$ be a synchronous cellular automaton or declared
local update. A run is the unbounded sequence $c_0,\, c_1 = T(c_0),\, c_2 = T(c_1), \ldots$ Fix a
regular reference background family $B_n$ with $B_{n+1} = T(B_n)$, or else with every mismatch between
$T(B_n)$ and $B_{n+1}$ explicitly recorded. A \emph{defect} at time $n$ is a finite, typed witness
that $c_n$ differs from $B_n$ on a declared local patch, and the \emph{defect census} is the typed
ledger
\[
  \delta_n \;=\; (D_n,\ \mathrm{create}_n,\ \mathrm{annihilate}_n,\ \mathrm{bind}_n,\ \mathrm{absorb}_n),
\]
in which every active defect at time $n$ is continued, annihilated, bound into a finite composite,
absorbed into a transporter, or explicitly marked unresolved at time $n+1$, and every new defect is
charged to a creation event.

A \emph{transporter} is a finite pattern record $R = (P, W, \tau, d, t_0, \mathrm{certificate})$ with
finite support window $W$, period $\tau > 0$, and displacement $d$. Its certificate is a checkable
recurrence-with-displacement equality: for every $k \ge 0$ in the certified regime and every site
$x \in W + k d$,
\[
  c_{\,t_0 + (k+1)\tau}(x + d) \;=\; c_{\,t_0 + k\tau}(x),
\]
with the same equality recorded after translating the finite support of $P$. This equality --- not
visual pattern recognition, and not holonomy language --- is what certifies motion.

Again the temporal quantifier carries the law: G9 quantifies over the whole tail of the run, asking
whether the defect census \emph{eventually} evacuates into recurrence-with-displacement records. Its
Foundations~IV cousin F19 (Object Persistence) has the opposite logical type: F19 starts from an
already-declared transport $\gamma : H_i \to H_j$ and checks whether an already-posited object
identity persists through it. In G9 the transporter is not declared at the start; it \emph{emerges}
over unbounded time from local defects. A one-dimensional rule over background ${.}$ with defect state
$X$ and transporter states $A,B$ separates the two: if an isolated $X$ becomes $AB$ in one step and
$AB$ thereafter shifts right forever, then at time $0$ there is no transport object for F19 to check
--- the record $(P{=}AB,\ \tau{=}1,\ d{=}{+}1)$ exists only after the evacuation event $X \to AB$ that
G9 audits. The row carries status tag \texttt{calibration-anchored schema}.

\begin{theorem}[\lawGNine]\label{thm:G9}
Let a run of an extended local system carry the setup data above.

\textbf{Positive schema.} Assume:
\begin{itemize}
  \item \textup{[H-G9-census]} the run carries an exact defect ledger, and there is a finite $N$ such
    that for all $n \ge N$ the unresolved defect measure is non-increasing after accounting for bound
    and absorbed defects: $|\mathrm{unresolved}_{n+1}| \le |\mathrm{unresolved}_n|$.
  \item \textup{[H-G9-absorb]} there are finite constants $N, T$ and a finite family of transporter
    records $\{R_j\}$ such that every defect still unresolved for $T$ steps after time $N$ is
    annihilated, bound into a finite non-moving composite, or assigned by the ledger to exactly one
    transporter record.
\end{itemize}
If the transporter-assigned persistent-defect set is nonempty and at least one assigned transporter
has displacement $d \neq 0$, then the run enters a certified transport regime: there is a finite time
$t_0$ and a finite nonempty family of transporter records covering every transporter-assigned
persistent defect after $t_0$, each record satisfying its recurrence-with-displacement equality on its
declared support. Persistent defects bound into stationary residue may coexist; the transport
certification rests on the moving subset. In the single-transporter or common-velocity case the
conclusion is eventual spatio-temporal periodicity modulo translation:
\[
  c_{\,t+\tau}\big|_{W+d} \;=\; \mathrm{shift}_d\!\left(c_t\big|_{W}\right)
  \qquad \text{for all certified } t \ge t_0.
\]
If instead all defects are eventually annihilated or resolved, with no persistent transporter-assigned
defect, the run is cleared or quiescent --- not transport-certified.

\textbf{Negative schema.} Separately, if the ledger exhibits unbounded unresolved creation ---
$\sup_n |\mathrm{unresolved}_n| = \infty$, or infinitely many creation events remain forever unmatched
by annihilation, binding, or absorption --- then the positive transport claim is blocked by the
proliferation ledger itself. Equivalently, such a ledger is incompatible with
\textup{[H-G9-census]}, so the positive schema's census obligation cannot be discharged; no
recurrence-with-displacement record may be inferred from finite local motion or holonomy language
alone.
\end{theorem}

\begin{proof}
G9 is a schema, so the proof spine is the required discharge pattern together with the two mechanized
fragments. The negative side is the census--proliferation incompatibility: if the census hypothesis
holds, some index $N$ starts a non-increasing unresolved tail, so every later unresolved count is
bounded by $|\mathrm{unresolved}_N|$; unbounded creation instantiated past that index would produce a
strictly larger later count --- a contradiction. The positive side is certificate packaging: the
absorption hypothesis exhausts every long-lived unresolved defect into annihilation, a stationary
composite, or exactly one transporter record; given a nonempty assigned subset with some $d \neq 0$,
each assigned record's finite recurrence equality is the transport witness, and stationary residue
does not disturb it. What the schema demands of any instance is precisely these three discharges:
eventually-controlled creation, exhaustion of persistent defects by absorption records, and the finite
recurrence certificate per transporter.

The rule-184 traffic automaton discharges the schema completely, as solved mathematics. The update
sends every adjacent pair $10$ to $01$, conserves particle number, and evolves finite periodic
configurations toward the known regimes: below density $\tfrac12$, free particle flow (blocked-car defects
evacuate); above $\tfrac12$, the dual free flow of holes (blocked-hole defects evacuate); at exactly
$\tfrac12$, the alternating critical pattern. The defect ledger is the count and arrangement of
blocked cars or holes; the transport record is the eventual periodic shift
\citep{Fuks1997,Nagel1996,ChowdhurySantenSchadschneider2000}.\footnote{Anchored to the substrate Lean
at \path|lean/SixBirdsFoundationsVI/Laws/G9DefectEvacuationTransport.lean|:
\path|census_excludes_unbounded_creation| for the negative schema,
\path|certified_transport_regime_of_assigned_nonzero| for the positive packaging. Status tag
\texttt{calibration-anchored schema}.}
\end{proof}

The certified statuses of a declared run are classified by priority --- finite-depth audits first,
then proliferation, then clearance, then transport, then bound residue:

\begin{center}
\footnotesize
\begin{tabularx}{\textwidth}{@{}lXX@{}}
\toprule
Case & Status & Witness type \\
\midrule
(a) Cleared / quiescent & All defects eventually annihilated or resolved; no persistent transporter or
bound residue. Not a transport claim. & Tail time $N$ with empty persistent set and closing ledger
records. \\
(b) Transport-certified & Nonempty persistent subset assigned to a transporter with $d \neq 0$;
stationary residue may coexist. & Records $(P, W, \tau, d, t_0)$ with checked recurrence equalities and
ledger assignments for residue. \\
(c) Defect-bound only & Persistent defects exist but every one has $d = 0$ or no transporter record. &
Bounded census tail plus stationary-composite certificates. \\
(d) Defect-proliferating & Creation unbounded or unmatched forever. & Ledger witness with
$\sup_n |\mathrm{unresolved}_n| = \infty$ or an infinite unmatched creation family. \\
(e) Undetermined at census budget & Only a finite-depth census or an unresolved tail obligation is
available. & Audit horizon $K$, defect budget, unresolved-event list, and an explicit no-claim beyond
$K$. \\
\bottomrule
\end{tabularx}
\end{center}

This is a certification partition for a declared run, not a decidability theorem for arbitrary
cellular automata. In interaction terms, G9 uses P5 packaging to identify defect composites and
transporter records and P6 ledger content to account for creation, annihilation, binding, and
absorption; the transport certificate is the finite recurrence equality, never an entropy-production
or arrow-of-time certificate. The lab record: G9-L1 ran exact rule 184 on a ring of length $2000$ at
five exact densities $\left(\tfrac{3}{10}, \tfrac{2}{5}, \tfrac12, \tfrac{3}{5},
\tfrac{7}{10}\right)$; all five runs reached the predicted evacuated defect census, and every
recurrence certificate was checked over $128$ consecutive tail windows, with $\tau = 1$ and
$d = \pm 1$ in the sub- and supercritical regimes respectively.

\subsection*{Layer instantiations}\label{layer-inst:G9}
\begin{description}
  \item[Rule 184 and traffic cellular automata --- calibration.] The solved density regimes above are
    the calibration instance discharging the census-to-transport mechanism
    \citep{Fuks1997,Nagel1996,ChowdhurySantenSchadschneider2000}.
  \item[Langton's ant --- conjectural.] The highway conjecture --- that every finite initial
    configuration reaches the $104$-step highway --- remains open; the standard ant is due to
    \citet{Langton1986}. What \emph{is} proved is trajectory unboundedness, due to
    \citet{BunimovichTroubetzkoy1992}; the rival ``Cohen--Kong theorem'' attribution circulating in
    popular accounts is flagged as incorrect by the specialist literature and is not used here. The
    unboundedness fragment is a nonclaim-scoped citation: G9's theorem and lab are carried by rule
    184, not by the ant.
  \item[Excitable-media automata --- structural.] Greenberg--Hastings-style systems produce discrete
    waves and spiral cores, but an instance requires an explicit defect census and transporter
    recurrence record for a specific model \citep{Greenberg1978,DurrettGriffeath}.
  \item[Cellular-automata glider ecology --- structural.] Life gliders and one-dimensional CA
    particles are direct recurrence-with-displacement objects once the period/displacement equality is
    checked on finite support; arbitrary debris-to-glider emergence remains model-specific
    \citep{Gardner1970,Pivato}.
  \item[Queue clearing and traffic shock waves --- structural.] Kinematic-wave and cell-transmission
    models supply moving jam and clearing-front analogues; an instance requires a discrete event
    ledger matching creation, annihilation, and absorption of queue defects
    \citep{LighthillWhitham1955,Richards1956,Daganzo1994}.
\end{description}

The Langton's-ant scope, the finite depth of recurrence certificates, and the reservation of
entropy-production vocabulary for genuine thermodynamic content are consolidated in
\Cref{sec:nonclaims}.

\subsection{\lawGTen}\label{sec:clusterC:g10}

Let $X$ be a state space and $D : X \to X$ a declared \emph{destructive} operator --- lossful in the
explicit sense that $D$ is not injective: some $x \neq y$ have $D(x) = D(y)$. Let $b : X \to B$ be a
declared boundary readout, $\rho : B \to B$ a declared boundary update, and $I : X \to \mathrm{Prop}$
an interior invariant. For an initial state $x_0$, the run is $x_{n+1} = D(x_n)$.

G10 asks whether the boundary relationship survives \emph{every} iterate of the lossy map, not merely
one application. Its Foundations~IV cousin F34 (Information Loss) characterizes when a \emph{single}
lossy step is legitimate --- when every later identification is invisible to the declared source
information. G10 points the other way: unbounded iteration of one lossful operator, with a positive
persistence conclusion, purchased not by reconstructing lost information but by a separately declared
invariant that regenerates the protection at every step. The row carries status tag \texttt{theorem}.

\begin{theorem}[\lawGTen]\label{thm:G10}
Assume:
\begin{itemize}
  \item \textup{[lossful]} $D$ is not injective;
  \item \textup{[H-G10-protect]} for every state $x$: $I(x) \implies b(D(x)) = \rho(b(x))$ --- the
    exact functional relationship, with no unstated target value;
  \item \textup{[H-G10-regen]} for every state $x$: $I(x) \implies I(D(x))$;
  \item \textup{[initial invariant]} $I(x_0)$.
\end{itemize}
Then for every $n \ge 0$: $I(x_n)$ holds, and the boundary obeys the iterated protection law
\[
  b(x_n) \;=\; \rho^{\,n}(b(x_0)),
  \qquad\text{equivalently}\qquad
  b(x_{n+1}) \;=\; \rho(b(x_n)) \ \text{ at every step.}
\]
The theorem is unconditional given the named hypotheses; it does not require $D$ to be invertible and
does not reconstruct lost interior information.
\end{theorem}

\begin{proof}
Induction on $n$. The base case is the assumed $I(x_0)$ together with $b(x_0) = \rho^{0}(b(x_0))$. For
the step, assume $I(x_n)$ and $b(x_n) = \rho^{n}(b(x_0))$. Regeneration gives $I(x_{n+1})$, and
protection gives
\[
  b(x_{n+1}) \;=\; b(D(x_n)) \;=\; \rho(b(x_n)) \;=\; \rho^{\,n+1}(b(x_0)).
\]
Both the invariant and the boundary law therefore hold for all $n$.\footnote{Anchored to the substrate
Lean at \path|lean/SixBirdsFoundationsVI/Laws/G10LossfulPersistence.lean|: the induction is
\path|boundary_persistence|; the worked positive witness is the checksum family
\path|checksum_lossful|, \path|checksum_protect|, \path|checksum_regen|,
\path|checksum_boundary_persistence|. Status tag \texttt{theorem}.}
\end{proof}

The hypotheses are cheap to state and expensive to own: the value of the law is the audit discipline
that separates a proved protection/regeneration pair from empirical persistence. The case enumeration
records that separation:

\begin{center}
\footnotesize
\begin{tabularx}{\textwidth}{@{}lXX@{}}
\toprule
Case & Status & Witness type \\
\midrule
(a) Certified-persistent & \textup{[H-G10-protect]}, \textup{[H-G10-regen]}, and $I(x_0)$ are proved
for the declared $D, b, \rho, I$. & The two hypothesis proofs plus the induction theorem. \\
(b) Collapsing & A finite certificate shows the proposed protecting structure is destroyed or the
state reaches a nilpotent terminal class. & Nilpotency or ranking certificate; Ducci at $k = 2^m$ is
the calibration. \\
(c) Persistent empirically at bounded depth & Checked only through finite depth $K$. & Exact
computation record through $K$, with no claim beyond. \\
(d) Unclassified & No regeneration proof, collapse proof, or bounded-depth certificate supplied. &
Missing-invariant or unresolved-obligation record. \\
\bottomrule
\end{tabularx}
\end{center}

The Ducci game calibrates the \emph{collapse} side literally. For $k = 2^m$ let $S$ be the cyclic
shift on $k$ coordinates; modulo $2$ the Ducci map $(x_1, \ldots, x_k) \mapsto (|x_1 - x_2|, \ldots,
|x_k - x_1|)$ is the linear operator $L = I + S$, since $|a - b| \equiv a + b \pmod 2$. In
characteristic $2$ the freshman's dream gives
\[
  L^{2^m} \;=\; (I + S)^{2^m} \;=\; I + S^{2^m} \;=\; I + I \;=\; 0
  \qquad \text{over } \mathbb{F}_2,
\]
because $S^k = I$. Each block of $k$ Ducci steps therefore deepens divisibility by $2$, while integer
Ducci iterates stay bounded after the first step; the tuple must collapse to zero. Any nonzero
boundary-persistence claim for all $2^m$-tuples therefore fails unless its invariant differs from the
nilpotently destroyed parity structure. In interaction terms, G10 concerns what survives repeated
P5-style compression and how P6 audit records certify the survival; it is not a reversibility claim
and not an entropy-production statement. The lab record: G10-L1 swept $k = 3, \ldots, 16$ with exact
integer arithmetic; the power-of-two lengths collapsed to zero (at steps $4$, $18$, $98$ for
$k = 4, 8, 16$), the GF(2) matrix checks confirmed $(I+S)^k = 0$ at exactly those lengths, and every
seeded non-power-of-two length entered a nonzero cycle.

\subsection*{Layer instantiations}\label{layer-inst:G10}
\begin{description}
  \item[Ducci games --- calibration, collapse side.] The mod-$2$ nilpotency argument above forces
    all-zero collapse for $k = 2^m$; non-powers of two are eventually periodic and may enter nonzero
    cycles \citep{ChamberlandThomas2004,Breuer2007a,Breuer2007b}.
  \item[Gilbreath / Proth--Gilbreath conjecture --- conjectural.] With $D$ the iterated absolute
    difference on prime rows and $b$ the leading entry, the target law is $\rho(2) = 1$, $\rho(1) = 1$
    with a small-gap/randomness interior invariant as the conjectural candidate. The conjecture is
    open; Odlyzko's computation supports the gap-smallness heuristic, and a precise random analogue
    exists \citep{Odlyzko1993,Chase2020}.
  \item[Finite-difference schemes and boundary-condition stability --- structural.] Repeated updates
    are lossy at unresolved scales while summation-by-parts and related boundary treatments can
    preserve declared boundary constraints; an instance requires the protection and regeneration
    estimates for a specific scheme \citep{Strikwerda,GustafssonKreissOliger,CoulombelFaye}.
  \item[Edge persistence under smoothing --- structural.] Scale-space smoothing destroys
    high-frequency detail while selected edge or zero-crossing structures may persist under declared
    conditions; the edge invariant and its regeneration remain per-operator obligations
    \citep{Canny1986,Witkin1983}.
  \item[Checksum cascades --- calibration, worked positive toy.] On $X = \{0,1\}^2 \times \{0,1\}$
    with $b((a,b),c) = c$, $\rho = \mathrm{id}$, invariant $I : c = a \oplus b$, and destructive map
    $D((a,b),c) = ((c,0),c)$, the map is non-injective, protection and regeneration hold, and the
    checksum boundary persists forever while the two-bit data block is genuinely lost. This is the
    mechanized positive witness; production checksum systems need their own algebraic invariants and
    remain structural \citep{Fletcher1982,StoneEtAl1998}.
\end{description}

The compatibility of G10 with the corpus's data-processing results, the open status of Gilbreath, and
the classical-novelty caveat for Ducci are consolidated in \Cref{sec:nonclaims}.

% =============================================================================
\section{Cluster D --- \clusterD}\label{sec:clusterD}

Everywhere else in the corpus, an obstruction is a liability: something to be priced, dissolved,
repaired, or confined. The three laws of this cluster record the opposite economy. In \lawGEleven{}
(G11), local constraints manufacture a finite defect against \emph{every} candidate global symmetry,
and the defects are the mechanism by which local rules force global aperiodic structure. In
\lawGTwelve{} (G12), a finite obstruction is never dissolved at all --- its survival inside a
homogeneous carrier is itself the certificate of a global impossibility. In \lawGThirteen{} (G13), a
generator that is thin in the strict arithmetic sense nevertheless saturates a thick target up to a
vanishing residual, with the genuine exceptions classified by type rather than explained away. G11
and G12 carry status tag \texttt{theorem}; G13 carries \texttt{schema}, its deep saturation inputs
being imported instance theorems rather than results reproved here.

\subsection{\lawGEleven}\label{sec:clusterD:g11}

Let $C$ be a finite local constraint system over configurations $x : \mathbb{Z}^d \to A$ with finite
alphabet $A$ --- equivalently, a subshift of finite type given by finitely many forbidden patterns;
the flagship case is $d = 2$ with Wang-tile or geometric matching rules. A nonzero $p \in
\mathbb{Z}^d$ is a \emph{candidate period} when $x(v + p) = x(v)$ for every $v$. A \emph{per-period
defect certificate} for $p$ is a finite region $R_p \subseteq \mathbb{Z}^d$ together with a proof
that the rules of $C$ restricted to $R_p$ are inconsistent with the equalities $x(v+p) = x(v)$
wherever both sides are represented --- the local-to-global gate: local rules do not imply global
anti-symmetry until the per-period obstruction is actually supplied. A \emph{hierarchy certificate}
is a forced composition structure: scales $L_k \to \infty$, finite block types $B_k$, and locally
readable marker data such that every admissible configuration decomposes uniquely into $k$-blocks at
every $k$, with $(k{+}1)$-blocks composed from $k$-blocks.

The Foundations~IV cousins are both single-object: F4 classifies one finite patch system by whether
one local family globalizes, and F40 diagnoses whether one declared symmetry descends through one
readout package. G11 is a $\forall p$ \emph{generative} forcing statement --- every candidate period
in an infinite family receives its own finite defect, and the hierarchy is the machine that generates
them. The row carries status tag \texttt{theorem}.

\begin{theorem}[\lawGEleven]\label{thm:G11}
Assume:
\begin{itemize}
  \item \textup{[H-G11-hierarchy]} every globally admissible $C$-configuration carries the declared
    hierarchy uniquely at every scale; the hierarchy is locally forced by the rules of $C$; and for
    every nonzero period vector $p$ some scale $k$ has forced marker or block structure that cannot
    be invariant under translation by $p$;
  \item \textup{[H-G11-nonempty]} at least one global configuration satisfies all local rules of $C$.
\end{itemize}
Then admissible configurations exist; for every nonzero $p$ the hierarchy supplies a scale whose
forced structure is incompatible with $p$-periodicity; because the forcing is local, the
incompatibility is witnessed on a finite region $R_p$, yielding a per-period defect certificate; and
therefore every admissible configuration is aperiodic.
\end{theorem}

\begin{proof}
The proof shape is Robinson-style hierarchical forcing. Local matching rules force small recognizable
blocks; blocks compose into larger blocks, and composition continues at unbounded scales. If a
configuration had period $p$, then at a scale much larger than $|p|$ the forced block-origin or
marker grid would have to be invariant under translation by $p$, contradicting unique composition.
The rules forcing that scale are local, so the contradiction lives on a finite patch --- the defect
certificate for $p$. Quantifying over all nonzero $p$ gives aperiodicity of every admissible
configuration.\footnote{Anchored to the substrate Lean at
\path|lean/SixBirdsFoundationsVI/Laws/G11GlobalAntiSymmetry.lean|:
\path|g11_global_anti_symmetry| (via \path|admissible_configurations_exist| and
\path|hierarchy_forces_aperiodic|), over \path|NonemptyAdmissible|, \path|Hierarchy|,
\path|ForcedBreaksPeriod|. Status tag \texttt{theorem}.}
\end{proof}

The theorem has an undecidability face that bounds what it can be asked to do. For an arbitrary
finite Wang tile set, deciding \textup{[H-G11-nonempty]} is the domino problem, undecidable by
Berger; so while a declared system is logically classified as exactly one of (a)
\emph{periodic-admissible} (some admissible periodic configuration, witnessed by a period plus a
finite fundamental domain), (b) \emph{aperiodicity-forced} (admissible configurations exist, all
aperiodic --- nonemptiness plus per-period defect certificates, with the hierarchy as the flagship
certificate generator), or (c) \emph{empty}, there is no general algorithm assigning an arbitrary
tile set to its case. G11 does not evade Berger's theorem; it types the certificates that settle
individual systems.

The classical constructions calibrate the mechanism across five decades. Berger's 1966 undecidability
proof produced the first aperiodic Wang tile set (commonly reported at $20426$ tiles)
\citep{Berger1966}; Robinson's 1971 six-prototile set carries the clearest forced square-hierarchy
proof --- the direct template for \textup{[H-G11-hierarchy]} \citep{Robinson1971}; Penrose's kite/dart
and rhomb tilings add matching rules with inflation/deflation hierarchy \citep{Penrose1974}; and the
2023 announcements give single-shape witnesses --- the hat (aperiodic monotile, reflections allowed)
and the spectre (chiral, translations and rotations only)
\citep{SmithEtAl2023Hat,SmithEtAl2023Spectre}. In interaction terms, G11 is P2 constraint gating, P4
hierarchy across forced scales, and P1 descent obstruction per candidate period --- the productive
obstruction is precisely local-rule incompatibility with periodic descent. The lab record: G11-L1 ran
exact SAT encodings of torus tilings; the monochromatic two-tile control set is periodic-admissible
($1 \times 1$ torus, SAT), while the Jeandel--Rao $11$-tile set produced UNSAT on every $p_1 \times
p_2$ torus with $1 \le p_1, p_2 \le 4$ --- a finite-period defect-certificate exhibit, not a decision
procedure.

\subsection*{Layer instantiations}\label{layer-inst:G11}
\begin{description}
  \item[Aperiodic tilings and Wang tile sets --- calibration.] Berger, Robinson, Penrose, and the
    hat/spectre monotiles are solved aperiodicity constructions on the hierarchy/per-period-defect
    mechanism \citep{Berger1966,Robinson1971,Penrose1974,SmithEtAl2023Hat,SmithEtAl2023Spectre}.
  \item[Quasicrystals --- structural/empirical bridge.] Long-range order without translational
    periodicity is experimentally real; a material is a G11 instance only after a precise local-rule
    model and hierarchy are supplied \citep{ShechtmanEtAl1984,Senechal1995}.
  \item[Symbolic dynamics and subshifts of finite type --- calibration/structural.]
    Multidimensional SFTs are the formal home of the classification; fixed-point aperiodic SFT
    constructions calibrate directly when a hierarchy proof is supplied
    \citep{LindMarcus1995,DurandRomashchenkoShen}.
  \item[Constraint satisfaction with forced non-repetition --- structural.] CSP encodings can force
    global asymmetry through local clauses; instances require nonempty global models and
    symmetry-defect certificates \citep{Dechter2003}.
  \item[Synchronization and coding patterns --- structural.] Synchronization strings prevent
    ambiguous repeated alignments by local labels; G11-like when the forbidden-period ambiguity has a
    finite certificate \citep{HaeuplerShahrasbi2017}.
\end{description}

The no-algorithm caveat, the specificity of the hierarchy-certificate pattern, and the
quasicrystal-model boundary are consolidated in \Cref{sec:nonclaims}.

\subsection{\lawGTwelve}\label{sec:clusterD:g12}

Let $X$ be a carrier with a graph relation $E \subseteq X \times X$ and let $\Gamma$ act transitively
on $X$ by automorphisms of $E$. In the flagship, $X = \mathbb{R}^2$, $E(x,y)$ means
$\lVert x - y \rVert_2 = 1$, and $\Gamma$ is the Euclidean isometry group. A proper $k$-coloring is
$c : X \to \{1, \ldots, k\}$ with $c(x) \neq c(y)$ whenever $E(x,y)$; for finite $W \subseteq X$,
$G[W]$ is the induced finite graph. The finite certificate is
\[
  \mathrm{cert}_k(W) \;=\;
  \bigl(\text{coordinates for } W,\ \text{edge list verified against } E,\
  \text{a proof that } G[W] \text{ is not } k\text{-colorable}\bigr),
\]
where the coordinate/edge part is not decoration: a non-$k$-colorable graph not actually realized at
unit distance in the plane is no Hadwiger--Nelson witness. \emph{Radiation} means exactly two things
--- orbit radiation (every $\gamma W$ carries the transported certificate, because $\Gamma$ preserves
$E$) and the compactness bridge (the de Bruijn--Erd\H{o}s principle, by which finite subgraph
chromatic numbers exhaust the global chromatic number). The Foundations~IV cousin F6 kills unpriced
needles against a declared package; G12 inverts the polarity --- the finite obstruction is not
dissolved, priced, or confined, and its permanent survival \emph{is} the global lower-bound
certificate. The row carries status tag \texttt{theorem}.

\begin{theorem}[\lawGTwelve]\label{thm:G12}
Assume
\begin{itemize}
  \item \textup{[H-G12-finite-witness]} some finite $W \subseteq X$ has a valid
    $\mathrm{cert}_k(W)$ proving $\chi(G[W]) > k$.
\end{itemize}
Then $X$ has no proper $k$-coloring: a global coloring would restrict to a proper $k$-coloring of
$G[W]$, contradicting the certificate; hence $\chi(X, E) > k$. Assume also
\begin{itemize}
  \item \textup{[H-G12-homogeneity]} $\Gamma$ acts transitively on $X$ preserving $E$.
\end{itemize}
Then every $\gamma W$ is an equally valid witness with the certificate transported by $\gamma$ ---
location-invariance of the witness, not a stronger bound. Assume further
\begin{itemize}
  \item \textup{[H-G12-compactness]} the metatheory includes the de Bruijn--Erd\H{o}s compactness
    principle for graph coloring --- a purchased set-theoretic commitment, needed for finite
    witnesses to \emph{exhaust} finite global chromatic bounds, and not needed merely to infer a
    lower bound from a displayed subgraph.
\end{itemize}
Then for finite chromatic bounds the global chromatic number of $(X, E)$ is the supremum of the
chromatic numbers of its finite induced subgraphs: finite witnesses are complete, not merely
sufficient, for finite lower-bound detection.
\end{theorem}

\begin{proof}
The lower bound is a restriction argument, not a metaphor: global colorings restrict to every finite
subgraph, so a finite non-$k$-colorable subgraph forbids a global $k$-coloring. Orbit radiation is
immediate from equivariance. The bridge has the usual de Bruijn--Erd\H{o}s shape: if every finite
subgraph is $k$-colorable, a compactness or ultrafilter-style choice argument assembles compatible
finite colorings into a global one \citep{BruijnErdos1951}.\footnote{Anchored to the substrate Lean
at \path|lean/SixBirdsFoundationsVI/Laws/G12FiniteWitnessRadiation.lean|:
\path|no_global_k_coloring| over \path|ProperKColoring| and \path|FiniteWitnessObstruction|;
\path|orbit_radiation| over \path|StructurePreserving|; and
\path|finite_witness_iff_no_global_coloring| gated by the named (not proved)
\path|CompactnessBridge|. Status tag \texttt{theorem}.}
\end{proof}

For Hadwiger--Nelson the current ZFC bracket is $5 \le \chi(\mathbb{R}^2) \le 7$, and the lower bound
is finite-witness radiation in exactly this sense. De~Grey's 2018 construction produced a
$1581$-vertex unit-distance graph that is not $4$-colorable, raising the bound from $4$ to $5$
\citep{DeGrey2018}; the Polymath16 program then organized witness reduction --- Heule reached
$553$ vertices by clausal proof minimization \citep{Heule2018}, and Parts reported a $509$-vertex,
$2442$-edge $5$-chromatic unit-distance graph, the smallest in this line \citep{Parts2020}. The
upper bound $7$ is the classical hexagonal seven-coloring. The exact value is open --- one of $5, 6,
7$ --- and the choice-sensitivity warning must be recorded rather than waved off: Shelah--Soifer-style
results show that plane and distance-graph coloring questions can be sensitive to the set-theoretic
background \citep{ShelahSoifer2003}, and measurable-coloring restrictions occupy the same axiom-ledger
row \citep{Falconer1981}. The finite lower bound itself does not depend on this subtlety once the
finite graph and its certificate are accepted.

A submitted claim classifies by audit priority: first fix the metatheory --- a claim relying on
global compactness, definability, or exact chromatic values whose truth depends on the background is
recorded as (c) \emph{choice-sensitive bridge} with an explicit axiom ledger; in a fixed metatheory a
proposed witness either (a) \emph{radiates} (valid certificate; all $\Gamma$-copies equivalent;
finite witnesses exhaustive under the bridge), or (b) is \emph{locally neutralizable} (invalid edge
list, unrealized coordinates, or an actual $k$-coloring), or (d) remains \emph{uncertified}. In
interaction terms, G12 is P1 obstruction plus P2 gating with a deliberately narrow P6 role ---
$\mathrm{cert}(W)$ is an audit record, not a debt ledger, and forcing more structure onto this law
would be artificial. The lab record: G12-L1 verified the vendored Heule/Polymath16 $826$-vertex graph
in exact algebraic arithmetic --- $826$ vertices, $4273$ edges checked against the unit-distance
relation, $2000$ sampled non-edges with zero unit-distance hits --- and confirmed non-$4$-colorability
by SAT ($3304$ variables, $22874$ clauses, UNSAT). The honesty note travels with the exhibit: this
$826$-vertex instance has no public DRAT certificate in Heule's \texttt{CNP-SAT} repository, so the
UNSAT determination is solver-verified rather than independently proof-checked (smaller graphs in
that repository --- $517$ to $803$ vertices --- do carry public DRAT certificates, which this lab
does not consume); the $509$-vertex Parts graph was not used because its raw coordinates were not
found publicly archived.

\subsection*{Layer instantiations}\label{layer-inst:G12}
\begin{description}
  \item[Hadwiger--Nelson --- calibration.] Finite $5$-chromatic unit-distance graphs certify
    $\chi(\mathbb{R}^2) > 4$; bracket $5$--$7$
    \citep{DeGrey2018,Heule2018,Parts2020,BruijnErdos1951}.
  \item[Finite CSP hardness gadgets --- structural.] A finite gadget certifies a bound for a problem
    family under a declared reduction --- witness radiation with a reduction map in place of $\Gamma$
    \citep{GareyJohnsonStockmeyer1976}.
  \item[Rigidity gadgets --- structural.] Finite frameworks certify rigidity properties under
    declared embedding rules \citep{Laman1970,AsimowRoth1978}.
  \item[Finite forbidden-minor witnesses --- structural.] An excluded minor certifies global
    nonmembership for every containing graph; the carrier is graph containment rather than a
    homogeneous motion group \citep{RobertsonSeymour2004}.
  \item[Mechanism-design impossibility gadgets --- structural.] Finite preference profiles witness
    impossibility for whole mechanism classes, with social-choice closure standing in for the
    $\Gamma$-orbit and bridge \citep{Gibbard1973,Satterthwaite1975}.
\end{description}

The open exact value, the purchased status of the compactness bridge, and the classical-novelty
caveats are consolidated in \Cref{sec:nonclaims}.

\subsection{\lawGThirteen}\label{sec:clusterD:g13}

Let $G$ be an algebraic group over $\mathbb{Q}$ with declared arithmetic lattice $G_{\mathbb{Z}}$,
let $\Lambda < G_{\mathbb{Z}}$ act on a discrete carrier $Y$, and let $O = \Lambda \cdot x_0$ be the
orbit of a base point, observed through a projection $\pi : O \to \mathbb{Z}$ against a target
$A \subseteq \mathbb{Z}_{\ge 0}$. The two adjectives in the law's name are strictly technical.
\emph{Thin} means Zariski dense with infinite index in the declared arithmetic lattice --- never
``vaguely sparse.'' \emph{Thick} means positive counting density --- for congruence targets, a finite
union of residue classes modulo some $Q_0$, with $d(A) = |R|/Q_0$ --- never ``visually large.'' An
\emph{expansion certificate} is a uniform spectral gap: a single $\varepsilon > 0$ bounding the
second eigenvalue of the congruence quotient graphs $\mathrm{Cay}(\Lambda/\Lambda(q), S_q)$ away from
$1$ for every admissible modulus $q$. The Foundations~IV cousin F4 tests one local datum against one
restriction map; G13 is an asymptotic orbit-coverage claim over an unbounded target family, and its
sibling relationship to G4 is real but mechanistically different --- G4 covers by target-dependent
certificates, G13 by orbit growth plus expansion. The row carries status tag \texttt{schema}.

\begin{theorem}[\lawGThirteen]\label{thm:G13}
This is a schema with imported mathematical payload. Assume:
\begin{itemize}
  \item \textup{[H-G13-thin]} $\Lambda$ is Zariski dense in the relevant algebraic group with
    $[G_{\mathbb{Z}} : \Lambda] = \infty$;
  \item \textup{[H-G13-expansion]} the congruence quotient family has a uniform spectral gap for the
    declared generators;
  \item \textup{[H-G13-congruence-audit]} $A$ is the finite union of residue classes modulo a
    declared $Q_0$ passing all declared local congruence obstructions --- explicitly a local audit
    only;
  \item \textup{[H-G13-imported-saturation]} an imported instance theorem proves
    $\#\bigl((A \setminus \pi(O)) \cap [1, N]\bigr) = O(N^{1 - \eta})$ for some $\eta > 0$, or a
    stronger finite-exception bound.
\end{itemize}
Then the projected orbit saturates the thick admissible target up to the declared exceptional-set
bound, and in particular the relative density of missed admissible targets vanishes:
\[
  \frac{\#\bigl((A \setminus \pi(O)) \cap [1, N]\bigr)}{\#\bigl(A \cap [1, N]\bigr)}
  \;\longrightarrow\; 0 .
\]
A finite-exception import gives full saturation beyond a threshold; a positive-density-only import
licenses no density-one claim. If additionally
\begin{itemize}
  \item \textup{[H-G13-reciprocity-record]} a declared infinite family $R_{\mathrm{rec}} \subseteq A$
    is proved missed by $\pi(O)$ by a reciprocity argument not expressible as exclusion from the
    finite congruence audit,
\end{itemize}
then the finite-congruence local-global claim is false for the instance --- without contradicting
density-one saturation when $R_{\mathrm{rec}}$ has density zero. The obstruction refines the residual
taxonomy; it does not overturn the saturation.
\end{theorem}

\begin{proof}
The schema's proof shape --- not reproduced here --- is: expansion gives effective equidistribution
of the thin orbit on congruence quotients; the audit fixes a positive-density admissible target; a
sieve or circle-method argument converts equidistribution and counting into the exceptional-set
bound. What the catalog itself proves is the elementary packaging: an eventually-sublinear missed
count against an eventually-positive-density admissible count has vanishing relative density, and a
reciprocity record is compatible with it.\footnote{Anchored to the substrate Lean at
\path|lean/SixBirdsFoundationsVI/Laws/G13ThinOrbitSaturation.lean|:
\path|relative_density_theorem| over \path|EventuallySublinear| and
\path|EventuallyPositiveDensity|, plus \path|reciprocity_compatibility_observation| over
\path|InfiniteFamily| and \path|ReciprocityCompatible|. Status tag \texttt{schema}.}
\end{proof}

The Apollonian flagship discharges every audit row with imported theorems. Descartes' circle theorem
constrains oriented curvature quadruples by
\[
  k_1^2 + k_2^2 + k_3^2 + k_4^2 \;=\; \tfrac{1}{2}\,(k_1 + k_2 + k_3 + k_4)^2,
\]
a quadratic form of signature $(3,1)$; replacing one curvature by the other root gives the four
Apollonian reflections, which generate the (thin) Apollonian group orbit of an integral root
quadruple. For the standard primitive packing $(-1, 2, 2, 3)$, reduction modulo $24$ yields the
admissible curvature classes $\{2, 3, 6, 11, 14, 15, 18, 23\}$ --- local density $8/24 = 1/3$.
Bourgain--Fuchs proved positive density for integral packings \citep{BourgainFuchs2010};
Bourgain--Kontorovich proved the density-one direction, with $O(N^{1-\eta})$ missed admissible
integers for effective $\eta > 0$ \citep{BourgainKontorovich2014Apollonian}; and Haag, Kertzer,
Rickards, and Stange proved in 2024 that the full local-global conjecture is nevertheless
\emph{false} for many packings: quadratic and quartic families pass the mod-$24$ audit but are missed
for reciprocity reasons --- obstructions of the thin group, not of its Zariski closure
\citep{HaagEtAl2024}. That is exactly the G13 distinction: congruence admissibility is a local audit,
and reciprocity is a separate residual type, not a disguised congruence failure.

The audit runs on two levels. System-level, a projected orbit is (a) \emph{saturated} (finite or
empty residual beyond a threshold), (b) \emph{density-one with exceptional set} (the
Bourgain--Kontorovich shape --- compatible with infinite zero-density residual families), (c)
\emph{positive-density only}, or (d) \emph{undetermined}. Target-level, after the local audit is
fixed, a specific integer or family is (i) congruence-obstructed, (ii) hit, (iii)
reciprocity-obstructed, or (iv) other-exceptional/uncertified --- with (i) and (iii) disjoint by
definition, since a reciprocity obstruction is recorded only after the congruence audit is passed. In
interaction terms, G13 is P2 gating plus P6 audit content, with P1 entering only through explicit
obstruction records; the fit is deliberately narrow, and treating spectral gap as a closure force
would overread the mathematics. The lab record: G13-L1 ran exact Descartes-reflection BFS for two
root quadruples to curvature bound $300000$, checking Descartes' identity on every generated
quadruple; both orbits stayed inside the expected mod-$24$ classes; the packing $(-6, 11, 14, 15)$
avoided every curvature in the published reciprocity-obstructed families $2n^2$, $3n^2$, $6n^2$
within the bound (zero hits --- the concrete case-(iii) exhibit), while the baseline
$(-1, 2, 2, 3)$'s residual missing set is recorded as case-(iv) uncertain within the bounded census,
not as a reciprocity obstruction.

\subsection*{Layer instantiations}\label{layer-inst:G13}
\begin{description}
  \item[Apollonian circle packings --- calibration, imported.] Thin Apollonian group, mod-$24$
    admissibility, spectral-gap input, density-one saturation, and reciprocity residuals
    \citep{BourgainFuchs2010,BourgainKontorovich2014Apollonian,HaagEtAl2024}.
  \item[Zaremba's conjecture --- calibration, imported.] Density-one results for denominators with
    bounded partial quotients --- the same thin-orbit saturation architecture
    \citep{BourgainKontorovich2014Zaremba}.
  \item[Affine sieve in thin orbits --- structural/calibration.] Almost-prime Pythagorean triples use
    the expansion-and-sieve architecture with almost-primality as the target
    \citep{KontorovichOh2012}.
  \item[Expander-based derandomized coverage --- structural.] Sparse generator walks cover test
    spaces under spectral-gap guarantees; the expansion certificate without the arithmetic
    reciprocity audit \citep{HooryLinialWigderson2006}.
  \item[Model-set diffraction supports --- structural, unverified for G13.] Model sets are sparse
    structured generators with rich observable spectra, but no thin arithmetic group or
    congruence/reciprocity audit is asserted \citep{BaakeGrimm2013}.
\end{description}

The imported status of all deep saturation inputs, the schema ceiling, and the strict senses of
``thin'' and ``thick'' are consolidated in \Cref{sec:nonclaims}.
% =============================================================================
\section{The extended currency taxonomy}\label{sec:currency-taxonomy}

The corpus taxonomy recognizes four structural roles for a currency --- a low-dimensional map
$u : Y \to \mathbb{R}^k$ attached to a carrier: \textbf{C1} \emph{constraint} (gates feasibility),
\textbf{C2} \emph{ledger} (adds along protocols), \textbf{C3} \emph{potential} (monotone Lyapunov
currency), and \textbf{C4} \emph{dual} (shadow price). Cluster A's three laws each carry a currency
that fits none of these roles exactly, and this section records the three extensions they force,
then proves the bridge from those currencies to the asymptotic-status vocabulary of the corpus's
audited-operational-realisability paper \citep{Tsiokos2026AOR}. Nothing in this section is a new
G-law: it is the taxonomy consequence of Cluster A, with one genuinely new (and elementary) theorem.

\begin{center}
\small
\begin{tabularx}{\textwidth}{@{}lllX@{}}
\toprule
Role & Name & Law anchor & Extension over C1--C4 \\
\midrule
C5 & amortized-potential & G3 & extends C3 by relaxing strict stepwise monotonicity to controlled
accounting under a potential \\
C6 & transfinite-escrow & G2 & extends C2 from additive finite ledgers to well-founded-valued
ledgers whose strict descent bars infinite nonterminal runs \\
C7 & solvency & G1 & a liveness-typed role: every nonterminal run eventually earns a native descent
certificate \\
\bottomrule
\end{tabularx}
\end{center}

The AOR paper defines four asymptotic statuses for a nonnegative sequence $c_0, c_1, c_2, \ldots$:
\emph{asymptotic\_budgeted} (some $B$ bounds $c_j$ for all $j \ge J_0$), \emph{summable}
($\sum_j c_j$ converges), \emph{contractive} (the defined ratios $c_{j+1}/c_j$ tend to a constant
strictly below $1$ on the nonzero tail), and \emph{eventual\_zero} ($c_j = 0$ for all $j \ge J_0$),
with the pure-sequence implication chain
\[
  \textit{eventual\_zero} \;\Longrightarrow\; \textit{contractive} \;\Longrightarrow\;
  \textit{summable} \;\Longrightarrow\; \textit{asymptotic\_budgeted}.
\]
The bridge below reuses these four names and their defining tail conditions --- which are pure facts
about nonnegative sequences --- for run-step-indexed sequences derived from G-law proofs. It does
\emph{not} identify G-law runs with AOR's colimit-refinement objects, and it transfers none of AOR's
categorical conclusions; the scope is exactly the sequences named.

\begin{theorem}[Cluster-A currency-to-status bridge]\label{thm:currency-bridge}
\textbf{(A) C7 solvency.} Fix the G1 data $(T, <, A)$ and a nonterminal point $x$, and let
$\mathrm{badtail}_j(x) \in \{0, 1\}$ indicate membership $x \in N_j$ in the $j$-th bad-tail
membrane. Then pointwise liveness at $x$ --- some $k \ge 1$ with $V_k(x) > 0$ --- holds if and only
if $\mathrm{badtail}_j(x)$ is \textit{eventual\_zero}. Globalized over nonterminal points, this is
the status-language reading of \Cref{thm:G1}.

\textbf{(B) C6 transfinite escrow.} Fix the G2 data and a run satisfying the stage-coherence audit
over a well-founded escrow carrier, and let $\mathrm{notdone}_k \in \{0, 1\}$ indicate that no stage
$i \le k$ has reached the terminal set. Then --- classically, or constructively once a concrete
terminal witness is extracted --- there is a finite stage $K$ reaching the terminal set, and
$\mathrm{notdone}$ is \textit{eventual\_zero}.

\textbf{(C) C5 amortized potential.} Fix a G3-style run with nonnegative actual costs $a_i$,
nonnegative amortized costs $\widehat{a}_i$, nonnegative potential $\Phi$, and write
$A_n = \sum_{i \le n} a_i$, $H_n = \sum_{i \le n} \widehat{a}_i$, $\Phi_0 = \Phi(s_0)$, so that
$A_n \le H_n + \Phi_0$ by \Cref{thm:G3}. Then the status of the amortized-cost sequence converts
into a total-actual-cost bound:
\begin{enumerate}
  \item \textit{asymptotic\_budgeted} ($\widehat{a}_i \le B$ for $i \ge J_0$) gives, for every
    $n \ge J_0$, $A_n \le \Phi_0 + H_{\mathrm{head}} + B\,(n - J_0 + 1)$ where
    $H_{\mathrm{head}} = \sum_{i < J_0} \widehat{a}_i$; in particular $A_n = O(n)$;
  \item \textit{summable} ($\sum_i \widehat{a}_i = S < \infty$) gives the uniform bound
    $A_n \le S + \Phi_0$ for every $n$: total actual cost over the infinite run is finite;
  \item \textit{contractive} reduces to the summable case by the geometric tail bound, with the
    explicit constant $A_n \le \Phi_0 + \sum_{i < J} \widehat{a}_i + \widehat{a}_J / (1 - q)$;
  \item \textit{eventual\_zero} ($\widehat{a}_i = 0$ for $i \ge J_0$) gives the fixed finite bound
    $A_n \le \Phi_0 + \sum_{i < J_0} \widehat{a}_i$ for every $n$.
\end{enumerate}
\end{theorem}

\begin{proof}
\textbf{(A)} The membranes fail downward: $N_k$ requires $V_j(x) \le 0$ for \emph{every}
$1 \le j \le k$, so once a descent certificate appears at some $j \le k$, membership fails in $N_{k'}$
for every $k' \ge k$. If liveness holds at $x$ with witness $k$, then $\mathrm{badtail}_{k'}(x) = 0$
for all $k' \ge k$; conversely, if the indicator is eventually zero, some finite membrane fails, and
unfolding that failure produces $1 \le j \le K$ with $V_j(x) > 0$. The global packaging is exactly
\Cref{thm:G1}.

\textbf{(B)} The indicator is monotone in the needed direction: once an accepted state appears by
stage $K$, it has appeared by every later stage. \Cref{thm:G2} rules out the infinite nonterminal
predicate; classically this yields a finite $K$ with the terminal set reached (and in a concrete
audited run the terminal stage is supplied directly), whence $\mathrm{notdone}_k = 0$ for all
$k \ge K$.

\textbf{(C)} Every case is the single inequality $A_n \le H_n + \Phi_0$ plus the standard fact about
nonnegative series matching the status: split into head and bounded tail
(\textit{asymptotic\_budgeted}); bound monotone partial sums by the finite limit
(\textit{summable}); dominate the tail by a geometric series (\textit{contractive}); observe the
tail sum is literally zero (\textit{eventual\_zero}).
\end{proof}

The Lean-traceability of \Cref{thm:currency-bridge} is deliberately not uniform, and the split
matters. Parts (A) and (B) are re-descriptions of the mechanized theorems
\path|part_a_reduction| and \path|no_infinite_nonterminal_run| --- no new mathematics, only the
demonstration that the AOR status vocabulary describes their conclusions exactly (Part (B)'s
extraction of a finite witness from the refuted infinite predicate is where the classical step, or
the concrete audited witness, enters). In Part (C), cases (i) and (iv) are literal readings of the
mechanized integer-valued theorems \path|telescoping_bound| and \path|uniform_actual_cost_bound|;
cases (ii) and (iii) are \emph{not} consequences of the landed Lean declaration --- over nonnegative
integers a summable tail is automatically \textit{eventual\_zero}, so the summable and contractive
cases have nontrivial content only over a scalar domain admitting genuine fractional decay. They are
paper-level mathematics: the same telescoping algebra applied over nonnegative rational or real
sequences, sound and elementary, but a scalar generalization of the mechanized theorem rather than an
instance of it. Case (ii) is nevertheless the main new content of this section: under a summable
amortized-cost tail, G3's finite-prefix safety bound becomes an infinite-run total-cost-finiteness
guarantee.

Concrete values make the four cases tangible (these are illustrations, not hypotheses of the
theorem): with $\widehat{a}_i = 2$, $a_i = 1$, $\Phi_0 = 5$, $n = 10$, the budgeted bound reads
$10 \le 25$; with $\widehat{a}_i = 2^{-i}$, $a_i = 2^{-(i+1)}$, $\Phi_0 = 3$, the summable bound
reads $A_n \le 1 + 3 = 4$ uniformly (partial actual sums approach $\tfrac{255}{512}$ at $n = 8$);
with $\widehat{a}_i = 3 \cdot 2^{-(i-1)}$, $a_i = \widehat{a}_i / 2$, $\Phi_0 = 4$, the contractive
constant gives $A_n \le 6 + 4 = 10$ (actual sum $\tfrac{189}{64}$ at $n = 6$); and with
$\widehat{a} = (5, 2, 0, 0, \ldots)$, $a = (3, 1, 0, 0, \ldots)$, $\Phi_0 = 1$, the eventual-zero
bound reads $4 \le 8$.

The scope boundaries of this bridge --- no functor into AOR's categorical machinery, no transfer of
its colimit conclusions, Cluster-A laws only, and the exact Lean-traceability split above --- are
consolidated with the other nonclaims in \Cref{sec:nonclaims}.

% =============================================================================
\section{Layer instantiations at a glance}\label{sec:instantiations-glance}

This section introduces no new claim. Its counterpart in \FoundIV{} maps catalog rows to the external
substrate papers that instantiate them; \FVI{} has no sibling substrate papers yet, so this section is
adapted rather than copied: it consolidates the sixty-five layer-instantiation rows already recorded
inside the thirteen law subsections into one navigational table. Each entry repeats the grade
vocabulary used throughout --- \emph{calibration} (a solved instance discharging the law's mechanism),
\emph{structural} (the shape matches but the law's named hypotheses have not been discharged for a
concrete model), \emph{conjectural} (the flagship question is open), and combinations thereof --- and
points back to the full description-list entry, where the citations live.

\begin{landscape}
\begin{table}[p]
\caption{All layer instantiations, by law. Grades abbreviated: cal.\ = calibration, str.\ =
structural, conj.\ = conjectural, imp.\ = imported instance theorem, unv.\ = unverified for this law.}
\label{tab:instantiations-glance}
\scriptsize
\begin{tabularx}{\linewidth}{@{}llXl@{}}
\toprule
Law & Cluster & Instantiations (grade) & Detail \\
\midrule
G1 \lawGOne & A & Collatz accelerated odd map (conj.\ flagship); binary-counter increment (cal.);
size-change/ranking termination (cal.); aliquot dynamics (conj.); switched and probabilistic
termination certificates (str.) & \Cref{sec:clusterA:g1} \\
G2 \lawGTwo & A & Goodstein sequences and hydra games (cal.); ranking-function termination (cal.);
term-rewrite termination by path orderings (cal.); proof-theoretic ordinal analysis (str.); AOR
cascade termination (cal., corpus tie) & \Cref{sec:clusterA:g2} \\
G3 \lawGThree & A & Binary-counter increment (cal.); dynamic-array doubling (cal.); splay trees
(cal.); passivity and storage functions (str.); energy-harvesting battery budgets (str.) &
\Cref{sec:clusterA:g3} \\
\addlinespace
G4 \lawGFour & B & Fibonacci--Sylvester Egyptian fractions (cal.); Erd\H{o}s--Straus (conj.\
flagship); Erd\H{o}s covering systems (cal.); symbolic integration by pattern families (str.);
CEGAR and proof automation (str.) & \Cref{sec:clusterB:g4} \\
G5 \lawGFive & B & Base-2 reverse-and-add (cal.); base-10 Lychrel candidates (conj.);
carry-propagation side channels (str.); finite-precision error fronts (str.); automatic sequences
and regular invariants (cal./str.) & \Cref{sec:clusterB:g5} \\
G6 \lawGSix & B & Recam\'an's sequence (conj.\ flagship); self-avoiding walks (str.); novelty search
and curiosity-driven exploration (str.); online algorithms with irrevocable commitments (str.);
greedy sequences in combinatorial number theory (str.) & \Cref{sec:clusterB:g6} \\
G7 \lawGSeven & B & Angel problem (cal.); pursuit-evasion and cop number (str./cal.); routing under
adversarial link failure (str.); percolation navigation above criticality (str.); robust motion
planning with dynamic obstacles (str.) & \Cref{sec:clusterB:g7} \\
\addlinespace
G8 \lawGEight & C & Abelian sandpiles and rotor-routing (cal.); quantitative confluence certificates
in rewriting (str.); distributed load balancing and gossip averaging (str.); neural avalanche models
(str., unv.); settlement netting and clearing vectors (str.) & \Cref{sec:clusterC:g8} \\
G9 \lawGNine & C & Rule 184 traffic cellular automata (cal.); Langton's ant (conj.);
excitable-media automata (str.); cellular-automata glider ecology (str.); queue clearing and traffic
shock waves (str.) & \Cref{sec:clusterC:g9} \\
G10 \lawGTen & C & Ducci games (cal., collapse side); Gilbreath / Proth--Gilbreath (conj.);
finite-difference schemes and boundary-condition stability (str.); edge persistence under smoothing
(str.); checksum cascades (cal., worked positive toy) & \Cref{sec:clusterC:g10} \\
\addlinespace
G11 \lawGEleven & D & Aperiodic tilings and Wang tile sets (cal.); quasicrystals (str./empirical
bridge); symbolic dynamics and subshifts of finite type (cal./str.); constraint satisfaction with
forced non-repetition (str.); synchronization and coding patterns (str.) & \Cref{sec:clusterD:g11} \\
G12 \lawGTwelve & D & Hadwiger--Nelson (cal.); finite CSP hardness gadgets (str.); rigidity gadgets
(str.); finite forbidden-minor witnesses (str.); mechanism-design impossibility gadgets (str.) &
\Cref{sec:clusterD:g12} \\
G13 \lawGThirteen & D & Apollonian circle packings (cal., imp.); Zaremba's conjecture (cal., imp.);
affine sieve in thin orbits (str./cal.); expander-based derandomized coverage (str.); model-set
diffraction supports (str., unv.) & \Cref{sec:clusterD:g13} \\
\bottomrule
\end{tabularx}
\end{table}
\end{landscape}

Two reading conventions carry over from the law sections. First, a \emph{calibration} grade attaches
to the mechanism the instance discharges, which may be narrower than the whole law: Ducci calibrates
G10's collapse side, not its positive persistence; the binary counter calibrates G1's currency
mechanism, not its ghost mechanism; and G6's two designed variants calibrate the shapes of its
schemas without upgrading the law's grade. Second, \emph{conjectural} flags an open flagship question
that the law organizes but does not resolve --- Collatz liveness for G1, Erd\H{o}s--Straus for G4,
Lychrel base-10 for G5, Recam\'an coverage for G6, the Langton's-ant highway for G9, Gilbreath for
G10 --- and the corresponding nonclaims are consolidated in \Cref{sec:nonclaims}.

% =============================================================================
\section{Deferred sketches and dropped candidates}\label{sec:sketches}

The catalog's size is a decision, not an accident: the mining pool contained roughly twice as many
eccentric problems as laws, and admission ran through the full protocol of \Cref{sec:methodology}
with the explicit understanding that killing a candidate is as valuable as admitting one. Two
candidates reached step 4 of the protocol and died there, in two different ways; both dispositions
are recorded here so the rejection reasoning remains auditable.

\paragraph{Hidden-Atlas Exactness --- dropped, absorbed.} The Somos-sequence integrality seed
(division at every step, provable integrality throughout) is fully absorbed by the corpus's existing
SAU pattern once its promoted object is named: the hidden cluster-algebra atlas under which
integrality is a shadow of Laurentness. No residual survives the comparison, so the candidate is not
a law; its correct home is a case-study paragraph in the SAU registry
(\Cref{sec:methodology} gives the full walkthrough). Nothing is deferred here --- the disposition is
final unless a genuinely temporal residual is someday exhibited.

\paragraph{Self-Observation Atomization --- parked, collided.} The look-and-say seed (Conway's
Cosmological Theorem: self-describing sequences decompose into a finite atomic chemistry) dies at
step 4 by collision rather than absorption: its residual --- a system whose rule is a theory of
itself --- lands in the territory where Foundations~V's carrier-side bounded-reflexivity apparatus
and this paper's theorist-side certification apparatus would have to be disentangled deliberately
\citep{FoundationsVCognition}. That disentangling is real research; the candidate is therefore
parked as a deferred sketch with an explicit pointer to Foundations~V, in exactly the way
\FoundIV{} parks its own dimension, smoothness, and nothingness candidates: deferred for lack of a
settled recognition interface, not dismissed as uninteresting. If the V/VI seam is ever resolved,
this is the first candidate to re-run through the protocol.

\paragraph{Scope trichotomy and future candidates.} Within this paper's own scope, every admitted
row lands in exactly one of the three statuses fixed in \Cref{sec:apparatus} --- \texttt{theorem},
\texttt{schema}, or \texttt{calibration-anchored schema} --- and the catalog asserts nothing about
rows it does not contain. The mining pool retains unmined problems, and lab work occasionally
suggests new seeds; the standing procedure is that a new candidate runs the full five-step protocol
and the admission gates \emph{before} any investment, is written up as a one-page candidate memo,
and enters the catalog only by explicit decision of the project owner. Thirteen laws is a signed-off
catalog, not a target to exceed. Likewise out of scope, explicitly: any substrate sequel applying
Cluster A to a full Collatz closure program, and any carrier-side (Foundations~V-style) treatment of
the G6/G7 obstruction landscapes.

% =============================================================================
\section{Nonclaims and caveats}\label{sec:nonclaims}

Nonclaims are load-bearing in this series: a law's value depends on what it declines to assert as
much as on what it proves. The individual law sections keep their prose focused on setup, statement,
and proof; this section consolidates the caveats, in four registers, and closes with the two
objections we expect and accept the burden of answering.

\paragraph{Open problems, left open.} This paper solves no open problem and reports no progress on
one. Specifically: the Collatz conjecture remains open (G1's Part~(b) hypotheses
[H-G1-ghost-convergence] and [H-G1-separation] are named obligations, and the $2^L - 1$ shadowing
family is evidence against fixed-depth strategies, not evidence about the conjecture); Tao's
almost-all result retains status \texttt{probabilistic} relative to G1's pointwise target and
discharges neither hypothesis. The Erd\H{o}s--Straus conjecture remains open (Mordell-style residue
coverage and the quadratic-residue obstruction are sub-results about identity families; density
evidence is never upgraded to pointwise closure). No base-10 Lychrel candidate is decided ---
$196$ remains open, and deep finite searches stay bounded evidence. Recam\'an coverage remains open
($852655$ is a bounded-census fact, not a permanent-hole proof). The Langton's-ant highway
conjecture remains open (trajectory unboundedness is proved, by Bunimovich--Troubetzkoy, and is
cited in a nonclaim-scoped role only). Gilbreath's conjecture remains open (the gap-smallness
reframing is heuristic). The exact chromatic number of the plane remains open (the bracket is
$5$--$7$). The Catalan--Dickson question for aliquot sequences remains open. Where a law names a
conjectural flagship, the law organizes the problem's structure; it never claims progress on the
problem itself.

\paragraph{No classical novelty.} The classical mathematics that anchors each law is imported, with
attribution, and no novelty is claimed for it: Dhar's abelian property and the Fey--Levine--Peres
least-action principle (G8); rule-184 traffic theory (G9); Ducci linear algebra (G10); Goodstein's
theorem, hydra termination, and the Kirby--Paris independence result (G2); classical amortized
analysis (G3); the Fibonacci--Sylvester algorithm and Erd\H{o}s--Straus partial results (G4); the
base-2 reverse-and-add result (G5); the angel-problem proof stack of
Conway--Berlekamp/M\'ath\'e/Kloster/Bowditch/G\'acs (G7); the aperiodic constructions of Berger,
Robinson, Penrose, and Smith--Myers--Kaplan--Goodman-Strauss (G11); de~Grey's graph, the
Polymath16/Heule/Parts reductions, and the de Bruijn--Erd\H{o}s theorem (G12); and the
Bourgain--Fuchs, Bourgain--Kontorovich, and Haag--Kertzer--Rickards--Stange theorems (G13, all
imported recognition sources). The contribution in every case is the layer-agnostic normal form, the
named hypotheses, the separating witnesses, and the audit discipline.

\paragraph{Conditionality where the register demands it.} Each law is exactly as conditional as the
corpus's nonclaim register requires, and several hypotheses are load-bearing by design. G8's abelian
compatibility is named and non-negotiable --- terminating confluence alone determines no canonical
counter ledger (the case-(b) witness), and Part~B's least action is the legal-versus-legal
comparison only. G9's transport is certified by recurrence-with-displacement records, never by
holonomy or visual motion; nothing in G8 or G9 is a drive, arrow, or entropy-production claim, and
that vocabulary stays reserved. G4 and G11 are conditional on their exhaustion and per-period
certificate data --- local admissibility never globalizes for free, and G11 supplies no algorithm
against Berger's undecidability. G5's no-go half is not an escape certificate; its confinement half
requires a verified target-free closed language. G6 and G7 are theorist-side throughout: the system
never certifies itself, and the adversary's strategy space is always declared --- no hidden-oracle
devil. G6's two rate hypotheses are proposed formalizations that later labs may refine, and its
designed-variant calibrations deliberately do not upgrade its grade. G7's thresholds are board- and
rule-relative; its finite-board lab is an exhibit, not infinite-board evidence. G10 does not
contradict the corpus's information-loss results and reconstructs nothing lost; its hypotheses must
be proved, not observed empirically. G12's compactness bridge is a purchased set-theoretic
commitment, recorded, and choice sensitivity in plane-coloring questions is a real warning; the
826-vertex lab witness is solver-verified UNSAT without an independent proof-checked certificate,
as its provenance record states. G13's status ceiling is \texttt{schema}; ``thin'' and ``thick''
carry only their strict technical senses, and congruence admissibility is local and incomplete ---
the 2024 reciprocity obstruction shows exactly that. Throughout Cluster A, promoted objects
(ledgers, escrows, completions, ghosts) are never literalized at the native layer; only boundary
certificates descend.

\paragraph{The currency bridge's scope.} \Cref{thm:currency-bridge} constructs no functor into
AOR's categorical colimit machinery and transfers none of its colimit conclusions; it reuses four
status names whose defining tail conditions are pure facts about nonnegative sequences. Its Parts
(A) and (B) are re-descriptions of proved theorems; in Part (C), the summable and contractive cases
are paper-level scalar generalizations of the integer-valued mechanized theorem, not corollaries of
it. The bridge is Cluster-A-scoped; no claim is made that every G-law currency bridges to the
ladder.

\paragraph{Two anticipated objections.} \emph{``These are known theorems relabeled.''} The rebadge
discipline of \Cref{sec:methodology} is the answer in mechanism: every law must differ from every
existing catalog row on a declared axis, the shared differentiating axis here is the temporal
quantifier that no static row possesses, and the protocol demonstrably kills candidates that fail
the test --- the two negative walkthroughs are the receipts. What is classical in each law is
imported and attributed; what is new is the normal form that makes thirteen scattered mechanisms
one vocabulary. \emph{``Eccentric problems are unserious.''} The detector logic says the opposite:
a problem with a tiny rule and a global obstruction is the cleanest available microscope for
exactly one structural mechanism at a time, and every cluster of this catalog is anchored by
solved, respected mathematics --- Kirby--Paris, Dhar and Fey--Levine--Peres, Berger through the
2023 monotiles, de~Grey and Polymath16, Bourgain--Kontorovich --- rather than by the folklore that
makes the problems eccentric.

% =============================================================================
\appendix
\section{Lean formalization}\label{app:formalization}

All thirteen laws are mechanized in Lean~4 in the paper's companion repository, under the namespace
\path|SixBirdsFoundationsVI.Laws|. The formalization is deliberately austere: every law file is
Mathlib-free with \emph{zero imports}, standalone against the Lean core, and the full tree builds
cleanly (\texttt{lake build}, currently $76$ jobs) with no \texttt{sorry}, no \texttt{admit}, and no
project-declared axioms anywhere in the thirteen law files --- a discipline enforced by a standing
per-file grep at every landing step.

\subsection{Wording discipline}\label{subsec:lean-wording-discipline}

Every claim in this appendix distinguishes ``proved in Lean'' from ``proved in prose and cited''
from ``imported without proof.'' One distinction is load-bearing enough to state as a rule: a law's
status tag (\texttt{theorem} / \texttt{schema} / \texttt{calibration-anchored schema}) grades the
\emph{mathematical claim}, while a declaration's axiom-dependency list prices the \emph{Lean proof's}
foundational cost, and the two axes are independent. G13 is \texttt{schema}-graded --- its deep
saturation content is imported --- yet its two mechanized declarations are nearly axiom-free; G1's
Part~(a) is an unconditional \texttt{theorem} yet its Lean proof uses \texttt{Classical.choice}.

\subsection{Structural difference from Foundations IV: no opacity-guard rows}%
\label{subsec:opacity-guard-pattern}

\FoundIV's formalization appendix devotes two subsections to opacity guards and narrowing
disclosures, because several of its rows are Lean wrappers around trust-base axioms delegating to
not-yet-formalized substrate papers. \FVI{} has no analogue of either subsection, and the difference
is worth stating plainly rather than hiding: every one of the thirteen laws here is a direct, fully
mechanized Lean theorem or schema in this repository. There are no substrate-delegated anchors, no
trust-base axiom wrappers, and no opacity guards; consequently there is also nothing to narrow or
parametrically disclose.

\subsection{Tagged-axiom register}\label{subsec:tagged-axiom-register}

Since the project declares no axioms of its own, the register below is the genuine Lean-kernel
axiom-dependency profile: which of the core's three foundational axioms (\texttt{propext},
\texttt{Classical.choice}, \texttt{Quot.sound}) each law's key theorems use. Every row was
independently re-verified by a fresh \texttt{\#print axioms} run per declaration at landing time.

\begin{table}[htbp]
\caption{Kernel axiom profile per law. ``Free'' = fully axiom-free.}
\label{tab:axiom-register}
\footnotesize
\begin{tabularx}{\textwidth}{@{}lXl@{}}
\toprule
Law & Key theorem(s) & Axiom profile \\
\midrule
G1 & \path|part_a_reduction|, \path|part_b_discharge| &
\texttt{propext}, \texttt{Classical.choice}, \texttt{Quot.sound} \\
G2 & \path|no_infinite_nonterminal_run|, \path|not_nonterminal_of_stage_coherence|,
\path|no_descending_chain_from_acc|, \path|wellFounded_no_descending_chain| & free \\
G3 & \path|telescoping_identity|, \path|telescoping_bound|, \path|amortized_sum_le_scale|,
\path|uniform_actual_cost_bound| & \texttt{propext}, \texttt{Quot.sound} \\
G4 & \path|positive_schema|, \path|conditional_biconditional|, \path|fixed_package_no_go| & free \\
G5 & \path|no_fixed_depth_future_sufficient|, \path|no_future_sufficient_depth|,
\path|closed_pattern_persists| & free \\
G5 & \path|iterate_add|, \path|closed_pattern_confinement| & \texttt{propext} \\
G6 & \path|covering_schema|, \path|hole_forming_schema| & free \\
G7 & \path|escape_never_stuck| & free \\
G7 & \path|play_succ_from_initial|, \path|no_infinite_nat_descent|,
\path|escape_confinement_mutually_exclusive| & \texttt{propext} \\
G7 & \path|confinement_forces_trap_with_fuel|, \path|confinement_forces_trap| &
\texttt{propext}, \texttt{Classical.choice}, \texttt{Quot.sound} \\
G8 & \path|partA_order_independent_and_odometer|, \path|least_action| &
\texttt{propext}, \texttt{Quot.sound} \\
G8 & \path|caseB_not_abelian|, \path|caseB_stable_run_from_S_ends_N|,
\path|caseD_nonconfluent_witness|, \path|caseD_terminates_from_S| & \texttt{propext} \\
G9 & \path|census_bound_from_start|, \path|census_tail_bound|,
\path|census_excludes_unbounded_creation| & \texttt{propext} \\
G9 & \path|certified_transport_regime_of_assigned_nonzero| & free \\
G10 & \path|boundary_persistence|, \path|checksum_protect|, \path|checksum_regen|,
\path|checksum_boundary_persistence| & \texttt{propext} \\
G10 & \path|checksum_lossful| & free \\
G11 & \path|admissible_configurations_exist|, \path|hierarchy_forces_aperiodic|,
\path|g11_global_anti_symmetry| & free \\
G12 & \path|no_global_k_coloring|, \path|orbit_radiation|,
\path|finite_witness_iff_no_global_coloring| & free \\
G13 & \path|relative_density_theorem| & \texttt{propext} \\
G13 & \path|reciprocity_compatibility_observation| & free \\
\bottomrule
\end{tabularx}
\end{table}

Exactly two laws use \texttt{Classical.choice}. G7's confinement theorems need classical case-splits
on an abstractly undecidable \path|Stuck| predicate --- reviewed and accepted as a genuine,
appropriately scoped necessity of the abstraction level, not a smuggled shortcut. G1's reduction and
discharge need classical contradiction to recover existential liveness from the negated universal
membrane condition. Every other proof is fully constructive modulo \texttt{propext} and
\texttt{Quot.sound}, which are Lean-core proof-irrelevance and quotient bookkeeping rather than
mathematically substantive choices.

\subsection{Per-row formalization coverage}\label{subsec:formalization-coverage}

\begin{table}[htbp]
\caption{Law-by-law coverage. Namespaces are
\texttt{SixBirdsFoundationsVI.Laws.}$\langle$entry$\rangle$.}
\label{tab:formalization-coverage}
\footnotesize
\begin{tabularx}{\textwidth}{@{}lXll@{}}
\toprule
Law & Paper name & Status tag & Lean namespace entry \\
\midrule
G1 & \lawGOne & \texttt{theorem} (Part a) / \texttt{schema} (Part b) & \path|G1HiddenAmortizedSolvency| \\
G2 & \lawGTwo & \texttt{theorem} & \path|G2TransfiniteEscrow| \\
G3 & \lawGThree & \texttt{theorem} & \path|G3AmortizedCurrency| \\
G4 & \lawGFour & \texttt{schema} & \path|G4MovingCoverExhaustion| \\
G5 & \lawGFive & \texttt{calibration-anchored schema} & \path|G5CarryHorizonConfinement| \\
G6 & \lawGSix & \texttt{schema} & \path|G6EndogenousNeedleGeneration| \\
G7 & \lawGSeven & \texttt{calibration-anchored schema} & \path|G7AdversarialMobilityConfinement| \\
G8 & \lawGEight & \texttt{theorem} & \path|G8OdometerAbelianization| \\
G9 & \lawGNine & \texttt{calibration-anchored schema} & \path|G9DefectEvacuationTransport| \\
G10 & \lawGTen & \texttt{theorem} & \path|G10LossfulPersistence| \\
G11 & \lawGEleven & \texttt{theorem} & \path|G11GlobalAntiSymmetry| \\
G12 & \lawGTwelve & \texttt{theorem} & \path|G12FiniteWitnessRadiation| \\
G13 & \lawGThirteen & \texttt{schema} & \path|G13ThinOrbitSaturation| \\
\bottomrule
\end{tabularx}
\end{table}

\subsection{Coverage summary by cluster}\label{subsec:lean-coverage-by-cluster}

\begin{center}
\small
\begin{tabularx}{\textwidth}{@{}lXccc@{}}
\toprule
Cluster & Laws & \texttt{theorem} & \texttt{schema} & \texttt{cal.-anchored} \\
\midrule
A --- \clusterA & G1, G2, G3 & G2, G3 + G1(a) & G1(b) & --- \\
B --- \clusterB & G4--G7 & --- & G4, G6 & G5, G7 \\
C --- \clusterC & G8, G9, G10 & G8, G10 & --- & G9 \\
D --- \clusterD & G11, G12, G13 & G11, G12 & G13 & --- \\
\midrule
Total & 13 & 6 full + G1 split & 4 full + G1 split & 3 \\
\bottomrule
\end{tabularx}
\end{center}

\subsection{Reviewer-cycle log}\label{subsec:reviewer-cycle-log}

Two review layers ran throughout the program. An \emph{internal} layer gated every artifact --- law
prose, Lean subsection, lab --- before anything built on it: manager verification (re-running builds,
re-deriving proofs, re-executing labs) plus an independent reviewer thread issuing ship/fix-first
verdicts. Once all thirteen laws were internally landed, a \emph{real external review cycle}
(requests v1 through v18, closing 2026-07-07/08, run by a separate external agent) provided the
genuinely adversarial pass --- and it found real things the internal layer had missed.

A sample of the internal layer's own catches, illustrative rather than exhaustive: G8's least-action
monotonicity hypothesis was initially vacuous (an early encoding discharged nothing) and was
corrected to a genuine exchange property before the Lean file was accepted; G7's devil strategy was
initially non-reactive --- depending only on the pre-turn state, not the angel's just-chosen move ---
a real fidelity gap against the prose turn order, fixed by threading the chosen position through the
strategy and legality data; and a forbidden sibling-series terminology string was found silently
propagated through seven gate files via a shared boilerplate sentence and swept in one pass.

The external cycle produced two genuine mathematical/Lean findings, both independently re-derived by
the manager against the actual source before any fix, and both closed with a fresh countermodel or
test proving the fix effective rather than merely compiling. First, G8 (request v1): the mechanized
\path|Sufficient| predicate required a legal run, so \path|least_action| proved minimality only
against other \emph{legal} sufficient scripts --- strictly narrower than the cited
Fey--Levine--Peres principle, which admits force-fired scripts with illegal intermediate states. The
resolution was honest scope narrowing: prose, gate records, and doc-comments were corrected to
describe the mechanized result precisely (\Cref{sec:clusterC:g8} carries that narrowed statement),
rather than attempting the strictly harder generalization, which would require unconditional rather
than legality-guarded commutativity. Second, G7 (request v7): \path|ConfinementStrategy| required
the devil's obstruction to decrease a rank but never required the obstruction to be \emph{legal},
admitting an illegal instant-trap countermodel --- a second, distinct devil-legality gap in the same
law. The fix added the legality field and threaded it through three theorems, and was verified by
constructing the reviewer's exact countermodel as a fresh Lean proof and confirming it is now
provably excluded. The remaining external findings were documentation and citation precision (stale
status wording in G1 and G6; two citation corrections in G13; a data-provenance correction in G12
recording the public certificate repository for smaller unit-distance graphs); G2, G3, G4, G5, G9,
G10, and G11 came back clean.

\subsection{Validator instructions}\label{subsec:validator-instructions}

This paper deliberately ships no build pipeline, lint script, or CI configuration; the validation
record is three items. (1)~\texttt{lake build} from the repository's \texttt{lean/} directory
compiles the full tree, including all thirteen law files, cleanly ($76$ jobs at the time of writing).
(2)~\texttt{pytest tests} from the repository's \texttt{lab/} directory runs the lab suite ($65$
tests at the last full run) covering the fast and structural checks of every lab exhibit cited in
this paper. (3)~The reviewer-cycle log above is the plain-language validation record: an internal
gate on every artifact, and external review requests v1--v18 closing every one of the thirteen laws.

\clearpage
\bibliographystyle{plainnat}
\bibliography{references}

\end{document}
